\documentclass[twoside,12pt]{article}
\usepackage{QSdocs}

\begin{document}

\noindent Docu on mex-functions and QSpace member functions\\
Auto-generated via QSdocs.m $\to$ QSdocs.tex $\to$ QSdocs.pdf\\
A. Weichselbaum \copyright{} \today

\tableofcontents

\newpage


\mysec[MEX routines (bin)]{MEX routines (bin)}


\pagebreak[3]


\mymex{NRGWilsonQS} % l=158/0

\begin{verbatim}
   Usage: [[NRG_DATA,] info] = ...
            NRGWilsonQS(H0, A0, Lambda, ff, FC, Z, [, gg, FL, OPTS])
 
   NRG run using (non-)Abelian symmetries for a Wilson chain of length L
   based on the standard iterative diagonalization prescription.
 
   The routine perfoms a quick check whether any of the inputs is complex.
   This includes the QSpaces H0, A0, FC, FL (if preseent), or the hopping
   amplitudes ff. If so, NRGWilsonQS runs in full complex arithmetic,
   thus returning complex AK and AD. The subsequent fdmNRG_QS also can
   deal with either case. All routines use double precision throughout.
 
      H0      impurity / starting Hamiltonian (rank-2 QSpace)
              H0 will be diagonalized upon initialization hence does not
              have to be diagonal on input. The resulting basis
              transformation is multiplied onto A0.
 
      A0      basis used for H0 in LRs order (rank-3 QSpace)
              where the local space s(igma) must already have the
              (equivalent) space of a Wilson site (note that A0 is used
              to obtain the matrix elements of FC at the first iteration.
 
      Lambda  used to rescale H_k by sqrt(Lambda); therefore this should be
              consistent with the Lambda used to obtain ff and gg below
 
      ff      strength of couplings along Wilson chain as they appear in
              the physical Hamiltonian, i.e. they must not be rescaled yet,
              as this is done automaticely based on Lambda! if ff is a vector
              of length L-1, the coupling applies for all nearest-neighbor
              hoppings via the FC operators in a diagonal fashion.     ^1)
              If ff is a matrix of dimension (L-1) x Nf, the diagonal
              hopping matrix elements can differ across the FC operators.
              If ff is a matrix of dimension (L-1) x (Nf*Nf), this also
              permits to encode cross-couplings FC(i)'*FC(j) in between
              nearest-neighbor shells.
 
              NB! The length L of the Wilson chain is deterimend by the
              length of ff, i.e. L=length(ff)+1 by including A0.
 
      FC      set of Nf annihilation operators used for coupling (QSpace)
              e.g. in case of all-abelian symmetries (default) this becomes
              H = F1'*ZF2 + F1*(ZF2)' + h.conj. (zflag=1, see below)
              with F1=F2=FC such that this operator alway appears as 
              a bilinear together with its hermitian conjugate.  ^2)
 
      Z       (diagonal) operator required for the calculation of the
              fermionic hopping matrix elements (QSpace). It is defined by
              ferminonic signs (-1)^ns required in the context as described
              for FC above, where ns is the number of fermionic particles
              for the local states s(igma).
   Options
 
      zflag   variation of how to apply the fermionic signs Z on the
              hopping terms in the Hamilonian determined by the operators
              FC (default: 1 e.g. as described for FC above in the context
              of all-abelian quantum numbers). Valid settings:
 
              0: ignore Z alltogether
              1: H = F1'*(Z*F2) + F1*(Z*F2)' + hermitian conjugate (default)
              2: intended for particle-hole-symmetric cases only:
                 H = (ZF1)'*(ZF2) at every second iteration only, WITHOUT
                 adding a further hermitian conjugate (already included!)
                 starting with NO Z-operators applied to A0.local (i.e.
                 same as for zflag==1), therefore starting with F1'*F1
                 at first iteration (i.e. using F1 to update operators,
                 and again F1 when building enlarged Hamiltonian.
                 For every other iteration, H = (Z*F1)'*(Z*F2) is applied
                 using Z with BOTH operators, consistent with particle-
                 hole symmetry!
              3: same as zflag=2, but Z is applied at every other iteration
                 already starting with the very first iteration.
 
      FL      set of Nl local operators along the chain (QSpaces [Nl])
 
      gg      strength of local operators FL along the chain similar to ff
              i.e. not rescaled! (double, (L-1) x Nl, not including H0)  ^1)
              Note that local operators can always be expressed as
              Hermitian operators, with real amplitudes gg. Hence gg is
              always expected to be real.
 
   ^1) NB! if the couplings ff (or also gg) are the same for all FC and FL,
       they may be represented by a single vector only.
 
   ^2) FC and FL are considered the same for every interation throughout
       the Wilson chain, yet there strength is modified by ff and gg,
       respectively.
 
   Further options / flags
 
    'Nkeep',. max number of states to keep at every NRG iteration (512)
    'NKEEP',. vector that specifies Nkeep individually for first few sites
 
    'Etrunc', use given energy threshold rather than Nkeep to trunctate
              unless Nkeep is reached (<=0 uses Nkeep only; default: 0).
    'ETRUNC', vector that specifies Etrunc individually for first few sites
    'ET1'     Etrunc for 1st with trunctation (default: 1.2*Etrunc)
              if ET1<0, then ET1=Etrunc is used.
 
    'FNfac',. global factor on energy scale, i.e.
              EScale(n) = Lambda^(-n/2) * 0.5*(Lambda+1) * FNfac.
              (by default, FNfac is chosen such that ff(N-1)/EScale(N)=1;
              see returned info.EScale below)
 
    'Estop',. stop dynamically once stable fixed point is reached, i.e.
              energies EK(1:Nkept/4) for iterations k and k+2 are the same
              to within given value for Estop (default: 0, i.e. not used).
    '-Estop'  same as 'Estop',1E-4
 
    'NEE',..  number of states to store for energy flow diagram (Nkeep*Ns)
    'fout',.  store data for every site into output file <fout>_##.mat
              if specified, the first RETURN argument is skipped (./NRG/NRG).
    'deps',.. restore degeneracy within eps (1E-12)
    'db',..   require minimum distance towards discarded energies (-1)
              (db<0 searches for maximum distance in vicinity of Nkeep)
    'dmax'..  maximum number of extra states to include for db (Nkeep/10)
              (mostly a safeguard in case db has been set too large).
    '-IdF'    build local Id from F-ops, fully ignoring A0 (NB! by default,
              sigma-space of A0 is considered a proper Wilson bath site!)
    'ionly'   info only (incl. Eflow) i.e. do not store iterative NRG data
    '-v'      verbose mode
 
   Output
 
    NRG_DATA  MPS data (only if 'fout' was not specified; otherwise this
              data is stored into specified file structure):
 
     .A[KD]   states K(ept) / D(iscarded) at current NRG iteration
              according to NRG truncation scheme (rank-3 QSpaces in LRs order)
     .H[KD]   eigenenergies corresponding to the eigenbasis of AK/AD;
              stored in diagonal format.
 
    info      info structure containing the following fields:
 
     .istr    info string
     .usage   (struct) info such as time stamp, etc.
     .HK      rescaled energy data for energy flow diagram (QSpace vector)
     .EE      collected eig(HK) data (purely numerical array)
     .E0      substracted ground state energy at energy scale of iteration
     .EK      [largest kept, lowest and largest discarded] energy for each iteration
     .EScale  NRG energy scale used along the Wilson chain (see FNfac above)
     .Lambda  Lambda used in the calculation (same as input)
     .ops     (struct) operator specific data for Wilson chain (FC,Z,ff,gg,etc.)
     .paras   (struct) technical NRG specific internal parameters
 
   Alternative usage: NRGWilsonQS('MATfilename')
   NRGWilsonQS('setup.mat' [, 'fout_tag'])
 
      This allows to read all input parameters from given setup.mat file
      (MatLab binary) which allows pure standalone run of NRG without
      requiring MatLab (except for the API for binary *.mat I/O).
 
      This way all variables and options MUST be defined by their names
      shown above; moreover, this usage must include a specification
      for fout_tag as all output data will be written there.
      fout_tag is either specified as 2nd argument above, else it should
      be defined in terms of the variable 'fout' in the parameter file
      setup.mat (default for fout_tag: './NRG/NRG').
 
   AW (C) Apr 2006 ; May 2010 ; Nov 2014
\end{verbatim}

\pagebreak[3]


\mymex{fdmNRG\_QS} % l=147/46

\begin{verbatim}
   Usage: [omega, A0 [,Info]] = ...
          fdmNRG(NRG_DATA [,Inrg], B, C, Z0, [,opts])
 
     This routine calculates arbitrary correlation function (CF)
     in the full density matrix Numerical Renormalization Group
     (fdm-NRG) framework at arbitrary temperature T for arbitrary
     abelian or non-abelian symmetries.
 
        A(w) = <B||C'>_w = FFT(<{B(t),C'}>).  (' equiv ^dagger)
 
     If B is specified as empty [], it will be automatically assumed
     equal to C, i.e. B=C as is the case for spectral functions;
     anti-commutator (default) or commutator will be chosen according
     to cflags below (see also zflags for Fermionic signs).
 
     This routine uses data from a previous standard NRG run with
     the result written as a matrix product state in QSpace format.
     It must contain the state space kept (AK) as well as the one
     thrown / discarded (AD) at every NRG iteration.
 
   See also: NRGWilsonQS
 
   Input:
 
     NRG_DATA   structure containing data of previous NRG run; if this
                contains path/file, then the data will be read from the
                files path/file_##.mat for every site (see NRGWilsonQS)
 
      .AK       MPS state of k(ept) space (regular NRG, QSpace)
      .AD       MPS state of t(runcated) NRG space (QSpace)
      .HK       effective NRG eigenspectra (kept, QSpace)
      .HD       effective NRG eigenspectra (truncated, QSpace)
 
     Inrg       NRG info structure (required if NRG_DATA contains
                structure; otherwise <NRG_DATA>_info.mat will be accessed.
                required elements within Inrg structure / data:
 
      .Lambda   NRG discretization parameter for the conduction band.
                Lambda is needed to undo the energy scaling permformed
                in obtaining the NRG data at the first place.
              
      .EScale   energy scale fore each Wilson shell
              
      .E0       (relative) ground state energy offsets in the same
                energy scales as H[KD]; this is needed for evaluating
                the correct weights of the density matrix
 
     B,C        operators acting on the impurity used to calculate the
                correlation between them as in A(w)=FFT(<{B(t),C'}>).
 
                If the operators B or C are specified by name, they
                will be read from workspace, and stored in NRG_DATA
                together with the other NRG MPS data (the latter is
                disabled by the flag `nostore')
 
                The operators in B and C are paired up for correlation
                functions. They are applied in the order of priority
                (whichever fits first): L, then s, then R onto A0
                as in <NRGDATA>_00 with index order LRs where L = left,
                R = right = combined, s = local state space sigma).
 
     Z0         Z0 operator for the space described by B and C
                required for fermionic setup where must represent
                Z0=(-1)^(number of particles); Z0 is directly
                contracted onto the B and C where 'zflags' is set.
 
   Options / Flags (default method: fDM)
 
    'fDM'       full-density-matrix approach to CFs (default)
    'NRho',k,   build density matrix from single iteration k
                (for testing purposes, only)
 
    'T',...     temperature at which correlation function is calculated;
                default: about an order of magnitude larger than the
                energy scale at the end of the Wilson chain
 
    'vflag',..  more detailed output with increasing value (default: 1)
    'partial'   store G1 and G2 separatly (e.g. for detailed balance
                or the calculation of expectation values)
    'nostore'   do not store operator elements of B or C with NRG_DATA
    'locRho'    if Rho needs to be calculated it is stored in RAM only
                and will be deleted after the program is finished
    'noRHO'     do not calculate Rho, eg. calculate partition function
                only (relevant with empty ops B and C only)
    'rhoNorm',. allows to specify the weights rhoNorm explicitely
                as input (intended for testing purposes only)
 
    'zflags'    on which operators in <B||C> to apply fermionic signs Z0
                ex. 'zflags', [ 1 0 ... ]  would apply Z0 to the first
                but not to the second operator (default: all 0 if Z0
                is empty, 1 otherwise).
 
    'nlog',...  number of bins per decade for log. discretization (256)
    'emin',...  minimum energy for omega binning (TN/1000)  ^1)
    'emax',...  maximum energy for omega binning (10.)      ^1)
 
   ^1) NB! first (last) bin contains all data below (above) emin (emax)
 
   Output
 
     om         omega corresponding to the energy binning [-emax,+emax]
     A0         raw data for spectral function (i.e. discrete in om)
 
     Info       info structure: specific fields are
       .ver     version / mode of code
       .T       actual temperature used
       .rho     FDM weight distribution for given temperature
       .RHO     reduced thermal density matrix for A0 (R-basis).
       .Om      Omega(T) = -T*log(Z), (grand)canonical potential.
 
       .A4      individual spectral data for each of the two
                contributions to the (anti)commutator (hence A4 has
                twice as many  columns as A0 above).
 
       .a4      (1st row) integrated spectral integrals over the two
                individual contributions to the (anti)commutator of
                given correlator (2nd row may be ignored; it refers
                to check w.r.t. detailed balance).
 
       .reA0    contains Re(G(omega->0)) for the calculated correlation
                functions. Importantly for finite temperatues, the
                correlators with cflags!=0, i.e. susceptibilities,
                acquire in addition to the Kramers Kronig transformation
                (principal value integral!) also a Matsubara correction
                for Ea==Eb, which would be skipped in the principal
                value integral; this correction [see Hanl and AW, PRB 2014]
                is added to the first term in the commutator.
 
       .symfac  IROP factors for given spectral functions.
 
   ALTERNATIVE USAGE
 
   fdmNRG('--mat','nrgdata', 'setup.mat' [,'B'], 'C')
 
     B and C are the same as above (B=C if B is not specified).
     This usage allows to read all input parameters from given MAT files
     (MatLab binaries); all variables and options MUST be defined in
     <setup.mat> by their names shown above; e.g. the operators B and
     C are specified by their names as defined in <setup.mat>.
 
     The output will be written / updated to the file set
     `<nrgdata>_##.mat' and `<nrgdata>_info.mat'.
 
     Output will be appended to setup.mat.
 
   See also NRGWilsonQS.
 
   AW (C) May 2006 ; May 2010 ; Nov 2014
\end{verbatim}

\pagebreak[3]


\mymex{fgrNRG} % l=152/22

\begin{verbatim}
   Usage: [[[om,] a0], Ifgr] = fgrNRG(NRG0, NRG2, Lambda, B [,C], [,opts])
 
      full density matrix (FDM) Numerical Renormalization Group (NRG)
      calculation for Fermi Golden Rule (FGR) rates, with extension
      to resonanant ineleastic x-ray scattering (RIXS).
 
   FGR mode (default)
 
      This computes the FGR spectral data (om,a0) witin the fdm-NRG
      framework at arbitrary temperature (default T=TN^+ with N the
      length of the Wilson chain). When Fourier tansformed, it corresponds
      to the time-dependent correlation function [see Ifgr.Itdm.at]
 
         a(t) = < B(t)C'> = tr[ rho_0 * B(t) C' ]
 
      Here B and C are local operators acting at or in close vicinity
      of the impurity. They are assumed to switch back and forth
      between two disconnected sectors of Hilbert space (e.g. by the
      presence of an (otherwise unaccounted for) valence core hole.
      The time evolution C(t) = exp(iH0*t) C exp(-iH*t) therefore becomes
      `mixed' with `initial' Hamiltonian H0 and final Hamiltonian,
      iteratively diagonalized for both, NRG0 and NRG2, respetively.
      The frequency correponds to om=E_final-E_initial.
 
      If C is not specified or empy, it is automatically assumed
      equal to B. This is the typical situation in FGR.
      The inital state is determined by the full density matrix (FDM)
      rho_0 of NRG0 at given temperature (default T ~ T_N ~ 0^+K).
      The dynamics occurs within NRG2.
 
   Input
 
      NRG*       (string pointer to) data of previous NRG runs;
                 if this contains <path/file>, then the data will be
                 read from the files <path/file>_##.mat files as
                 generated by NRGWilsonQS.
 
      Lambda     NRG discretization parameter used in NRG*.
                 While EScale and E0 will be read from NRG0_info,
                 Lambda is needed for simple estimates of energy scales.
 
      B,C        local operators, e.g. acting on the impurity.  *)
                 If operators B or C are specified by name, they will be
                 read from workspace and stored in NRG2_info
                 (this is disabled by `nostore').
 
      NB! Note that exchanging NRG0 <> NRG2 and B<>C corresponds
      to changing from emission to absorption spectra or vice versa.
 
   Options / Flags
 
     'Eoffset',. additional energy offset which affects the omega
                 binning (om-binning) only. Eoffset acts *on top* of
                 E0_final-E0_initial, which ensures that by default,
                 (Eoffset=0) omega=0 corresponds to minimal excitation
                 energy for a transition from groundstate in NRG0
                 to groundstate in NRG2 (with NRG2 considered (much)
                 higher in energy in x-ray absorption).
 
     'T',...     temperature at which correlation function is calculated (0.)
     'NRho', k   builds density matrix from iteration k only (cf. `DM-NRG')
 
     'emin',...  min. energy for omega binning (TN/1000)  **)
     'emax',...  max. energy for omega binning (10.)
     'nlog',...  number of bins per decade for om-binning (256)
 
     'nostore'   do not store operator matrix elements of C[12] with NRG data
     'locRho'    if Rho needs to be calculated it is stored in RAM only
                 and will be deleted after the program is finished
 
     'sigma',..  applies broadening in frequency domain (log.gauss)
                 NB! either alpha or sigma can be specified but not both.
                 default sigma=0.3 (with alpha ignored)
     'alpha',..  exponential decay of energies exp(-alpha dE) [0.]
                 this applies alpha in time domain (F. Anders 2005)
 
     'partial'   return partial data from every NRG iteration
     'raw'       raw data only, i.e. no smoothing of data
     'RAW'       calculates time dependent data from exact frequency data
                 i.e. without binning of energies on log. scale
     'disp'      more detaild output (0)
 
     'Z0',..     Z0 operator for first site (impurity) required for
                 fermionic setup where Z0=(-1)^(number of particles)
     'zflags',.. on which operators in B.*C to apply fermionic signs Z0
                 ex. 'zflags', [ 1 0 ... ]  would apply Z0 to the first
                 but not to the second operator  (default: all 1s,
                 as typically the case for absorption/emission setups.
 
   RIXS mode
 
      With the flag '--rixs', this routine computes disrete, i.e., binned,
      RIXS (resonant ineleastic x-ray scattering) energy loss spetra
      for given incomming photon energy E_inc.
      This uses the Kramers-Heisenberg higher-order generalization of
      Fermi golden rule.
 
      Note on semantics: within RIXS, i(nitial) and f(inal) refer to the
      *same* Hamiltonian, wheras it is the intermediate Hamiltonian takes
      the role of the former H2:=H_final. With
 
          V_if := <i| B 1/[E_inc^* + E_i - H2] C' |f>
 
      the RIXS specral data would correspond, when Fourier transformed
      into time-domain, to
 
          R(t) = < V(t)*V'> = tr[ rho_0 * V(t) V' ]
 
      thus recovering the same final form as for absorption/emission
      spectra above. The RIXS mode requires the additional parameter
 
      'E_inc',.. incoming photon energy E_inc = complex(omega_in, gamma);
                 this must already also include the broadening gamma
                 related to finite life-time of the core-hole due to other
                 unaccounted for relaxation processes such as Auger etc.
 
      The returned spectra correspond to the RIXS intensity,
      i.e., spectral intensity a0 vs. energy loss om.
 
   Output
 
      om, a0     binned spectral data (according to emin,emax,nlog)
      Ifgr       info structure
 
   Notes / comments
 
    *) Alternatively, if B and C have composite structure that oct on
       multiple sites/indices simultaneously, the operators B and C can
       also be specified as cell arrays of QSpace vectors.
       A given QSpace vector then, e.g., B{i}, is interpreted as the
       sequence of operators that act like a tensor product on indices
       L00,s00,s01,... where L00 and s00 are the `left' and local (`sigma)
       indices of A0, respectively, and s01, etc. the local indices
       of A01, etc. In this usage, Z0 must be already properly included
       in B and C upon input, i.e. Z0 and zflags cannot be used.
 
   **) The bin a0(end) collects all data with omega >= om(end) = emax.
       For emin, the smallest bin a0(1) (given the logarthmic grid!)
       contains all data where omega is exactly zero (e.g. for diagonal
       contributions in Lehmann representation). The next smallest bin
       a0(2) then contains all data for 0 < omega <= om(2).
  
   Alternative usage:
   fgrNRG('nrgdata0','nrgdata2','setup.mat','B','C')
 
       This reads all input parameters and options from given setup.mat
       file. Hence the operators B and C must be specified by their name.
       Output will be written / updated to NRG0 and NRG2,
       with results also appended to setup.mat.
 
   See also NRGWilsonQS, fdmNRG_QS, tdmNRG
 
   AW (C) 2008 ; RIXS 2019
\end{verbatim}

\pagebreak[3]


\mymex{tdmNRG} % l=100/3

\begin{verbatim}
   Usage: [cc [,Info]] = ...
   tdmNRG(NRG0, NRG2, Lambda, C1 [,C2], tt, [,opts])
 
      time-dependent numerical renormalization group (NRG)
      using the full thermal density matrix of the initial system
      (FDM), rho_0, as starting configuration at time t=0.
 
   This routine calculates the time evolution of local operators
   C1(t)*C2 in fDM-NRG framework at given temperature T.
 
      cc(t) = < C1(t)C2'>_0 = trace( rho_0 * C1(t) * C2')
 
   With option -fgr the initial density matrix is modified
   by the operator C2 (e.g. creation / deletion of a particle
   as in absorption / emission processes) with the
   resulting modified time evolution being
 
      cc(t) = trace( (C2' rho_0 C2) * C1(t) )
 
   Hence the operator C2 appears twice. Note, however, that
   the projected density matrix C2' rho_0 C2 is no longer
   normalized.
 
   This routine uses data from two previous NRG runs with
   the result written as matrix product state in QSpace format.
   It must contain the state space kept (AK) as well as the one
   discarded (AD) at every NRG iteration.
 
   The inital state is determined by the density matrix rho_0 of
   NRG0 at given temperature T. The time evolution is determined
   within NRG2.
 
   See also: NRGWilsonQS, dmNRG, fgrNRG
 
   Input:
 
     NRG[02]    specifies previously calclulated NRG data; if this
                contains path/file, then the data will be read from the
                files path/file_##.mat for every site (see NRGWilsonQS)
 
        .AK     MPS state of k(ept) space (regular NRG, QSpace)
        .AD     MPS state of t(runcated) NRG space (QSpace)
        .HK     effective NRG eigenspectra (kept, QSpace)
        .HD     effective NRG eigenspectra (truncated, QSpace)
 
     Lambda     NRG discretization parameter for the conduction band.
                Lambda is needed to undo the energy scaling permformed
                in obtaining the NRG data at the first place.
 
     C[12]      local operators, e.g. acting on the impurity
 
     If C2 is not specified or specified as empty [], it will be
     automatically assumed equal to the identity operator.
 
     NB! If operators C[12] are specified by name, they will be
     read from workspace and stored in NRG_DATA together with
     the other NRG data (can be disabled by flag `nostore').
 
   Options / Flags (default method: tDM)
 
    'alpha',..  exponential decay of energies exp(-alpha dE) [0.]
                this applies alpha in time domain (F. Anders 2005)
    'sigma',..  applies broadening in frequency domain (log. gauss)
                NB! either alpha or sigma can be specified but not both.
                default: sigma=0.3 with alpha ignored.
 
    'NRho', k   builds density matrix from iteration k only.
    'T',...     temperature at which correlation function is calculated.
                If not specified, the temperature T is given by the
                energy scale towards the end of the Wilson chain.
 
    '-fgr'      run in FGR modus applying C2 twice (see above).
 
    'nlog',...  number of bins per decade for logarithmic discretization (256)
    'emin',...  minimum energy for omega binning (TN/1000)
    'emax',...  maximum energy for omega binning (10.)
 
    'nostore'   do not store operator elements of C[12] with NRG_DATA
    'locRho'    if Rho needs to be calculated it is stored in RAM only
                and will be deleted after the program is finished
    'partial'   return partial data from every NRG iteration
    'raw'       raw data only, i.e. no smoothing of data
    'RAW'       calculates time dependent data from exact frequency data
                i.e. without binning of energies on log. scale
    'disp'      more detaild output (0)
 
     Z0         Z0 operator for first site (impurity) required for
                fermionic setup where Z0=(-1)^(number of particles)
    'zflags'    on which operators C[12] to apply fermionic signs Z0
                ex. 'zflags', [ 1 0 ... ]  would apply Z0 to the first
                but not to the second operator  (default: none).
 
   NB! The first (last) bin contains all data below (above) emin (emax)
 
   Output
 
     cc         expectation value at times specified in tt
     Info       info structure
 
   See also NRGWilsonQS, fdmNRG_QS, fgrNRG
   AWb (C) Jan 2007
\end{verbatim}

\pagebreak[3]


\mymex{compactQS} % l=30/50

\begin{verbatim}
   Usage: S=compactQS([opts],'qtype',Q1,Q2,Q,A [,perm]);
 
      reduce matrix elements of the rank-3 IROP.
      Here A is still specified in full (dense) tensor format
      with matrix elements given by
 
         A(i1,i2,io) := <i1|A(io)|i2>
 
      NB! this assumes the index (1,2,op) specified on the lhs!
      The symmetry types are specified by qtype. Moreover,
      all states i1 and i2 must be already cast into symmetry
      eigenstate grouped into multiplets (see getSymStates.cc)
      with q- and z-labels both present and interleaved in
      each input state label sets Q*.
 
      The state spaces are combined in the order (Q2+Q =>Q1),
      thus using the Clebsch Gordan coefficients (Q1,Q|Q2).
      This is a natural order in that, Q=-1 annihilates a
      particle on a |ket> state (the annihilation operator
      on a ket state increases charge, i.e. Q would be +1,
      in case Q1 and Q had been combined into Q!).
      This therefore represents a natural index order.
 
   Options
 
      perm  permutation to be applied to {Q1,Q2,Q} and A
            prior to combining first two indizes into third.
 
      '-h'  display this usage and exit
 
   (C) AW : Apr 2010 ; Jul 2011 ; Oct 2014
\end{verbatim}

\pagebreak[3]


\mymex{contractQS} % l=138/27

\begin{verbatim}
  Usage: contractQS(A,..B,..)
 
      performs pairwise contraction of tensors
      e.g. A and B in QSpace format while also taking care of the
      underlying Clebsch Gordan coefficient spaces (if present).
 
      Each QSpace can be used as is, or as its 'conjugate' where the
      conjugate of a QSpace A, i.e. conj(A) is defined as the QSpace
 
      a) with all `arrows' reversed
      b) keeping the SAME qlabels [for reversing individual
         arrows, see getIdentitQS(..,'-0')]
      c) and complex conjugation of all data{*}
         if applicable, i.e. a given QSpace is complex
 
   Note that because of (a), the specification of conjugation flags
   (`conj-flags') is also important for QSpaces with all-real matrix
   elements.
 
   Usage #1: S=contractQS(A, ica, B, icb [, perm, OPTS ]);
 
      Plain contraction of a single pair of QSpaces, A and B,
      with respect to given explicitly specified  sets of
      contraction indices ica and icb, which can be specified
 
        - numerically (e.g. [1 2]),
        - or as strings (e.g. '1,2', or '1 2')
        - or as compact strings (e.g. '1,2', or '12')
 
      The last 'compact format' is only possible / unique,
      of course, if the tensors A and B do not have more than
      9 legs (which basically never occurs), such that the
      contraction indices reamin in the single digits
      (this can be further relaxed, though, by extending the
      digital range to letter, i.e. using a-z after 0-9).
 
      The recommended way to specify conj-flag with usage #1
      is together with the contraction indices in string notation!
      For example,
 
        contractQS(A,[1 3],B,[1 3],'conjA')  is equivalent to
        contractQS(A,'1,3;*',B,[1 3])        is equivalent to
        contractQS(A,'13*',B,'13')
 
      (Deprecated) options specific to usage #1:
 
        'conjA'  use complex conjugate of A for contraction
        'conjB'  use complex conjugate of B for contraction
 
   Usage #2: S=contractQS({A,{B,C}},... [, perm, OPTS ]);
 
      Generalized 'cell-contraction' of tensors: when encountering
      a cell, the content of that cell is contracted first, before
      using its results. This allows the specification of an entire
      patter of pairwise contractions based on nested structures
      where the lowest-level contractions are performed first.
 
      Cell contractions are furthermore based on QSpace 'itags'
      i.e. string labels for indices with up to 8chars, and which
      are specified in QSpace.info.itags. This offers automated
      contraction ('auto-contraction') of pairs of tensors solely
      based on matching itags! Uusage #2 therefore does not (also)
      allow explicit specification of contraction indices as in
      usage #1.  Therefore itags (plus conj-flags) must be unique.
 
      Itags must always also contain individual conjugate flags
      (this represent the bare minimum that must be specified with
      each QSpacein v3): the conjugate flags on individual indices
      (legs) of a tensor determine  whether that index (leg) is
      in- or out-going, with the convention that
 
          all out-going indices have a trailing * in their itags!
 
      For example, an A-tensor with L(eft), R(ight), s(=local)
      indices may have itags A.info.itags={'L','R*',s'}
      assuming (L,R,s) index order.
 
      In usage #2, for every operator additional optional strings
      can be specified, appearing right after the affected tensor
      e.g. QSpace A:
 
        A,'!ij'  do not contract indices specified by ij (in compact
                 format) despite they share common matching itags.
        A,'*'    apply overall (complex) conjugation on given input
                 tensor A (see early comments above)
        A,'!ij*' both of the above (with * always trailing).
 
        X,'-op:<tag>[:<opl=op>]' specify itags for given (e.g. local) operator.
 
           The last option considers X an operator, and hence assumes
           operator itags '<tag>;<tag>*[;opl*]' for QSpace X;
           the default operator label is `op', but may be changed
           by specifying a trailing ':<opl>' as indicated above.
           As a safeguard, this issues a warning, if existing itags are
           overwritten. This is relevant e.g. for local operators that
           are applied to a very specific site with associated itag.
 
           [11/24/2018] the specified <tag> may now also represent
           a regular expression (regexp), recognized by non-alphabetic
           special characters, while ignoring conj-flags (without
           special characters, the specified <tag> is taken as is!).
           This usage then searches for a matching itag in the paired
           up QSpace (cell) in the contraction.
 
           Ex. Consider A some QSpace with a single local index
           that starts with `s', e.g., like 's010' for site 10;
           then contract(A,Xloc,'-op:^s') will autocontract
           the local operator Xloc to the correct local index in A
           (here the regex `^s' indicates `starts with s');
           An alternative operator itag may still be specified
           by adding a trailing ':opl' as indicated earlier.
 
      An adaptation of usage #2 also be used for a plain sequential
      contractions where
          S=contractQS(A,B,C,... [, perm, OPTS ]); is equivalent to
          S=contractQS({A,{B, {C,...}}}, [,perm,OPTS ]),
      i.e. the sequential contractions are started from the end
      onwards to the beginning of the set.
 
   The remaining trailing OPTS are
 
      perm  permutation to be applied to the final object;
            NB! [06/02/2019] this permutation can be shorter
            than the rank of the resulting QSpace; in this case
            it only affects the leading range of indices.
 
      '-v'  verbose mode that shows level of cell contraction
            together with actual contractions performed.
 
   Mixed usage of #2 and #1 is not possible.
 
   Note that mex files do not allow to return class objects,
   hence S is returned as a QSpace structure. To get a class
   object, use QSpace(contractQS(...)), or equivalently,
   if contract() is properly defined as a wrapper routine
   within matlab's Class/@QSpace (see MPS Pack), having a QSpace
   input A, this may be shortended to contract(A,...).
 
   AW (C) May10 ; Aug12 ; Dec14 ; Sep16
\end{verbatim}

\pagebreak[3]


\mymex{diagQS} % l=4/53

\begin{verbatim}
   Usage: dd=diagQS(A)
 
      Get diagonal data of given QSpace (for rank-2 QSpaces only).
 
   Wb,Aug16,16
\end{verbatim}

\pagebreak[3]


\mymex{eigQS} % l=39/4

\begin{verbatim}
   Usage: [E [,I]]= eigQS(H [,opts])
 
       Obtain eigenspectrum and eigenbasis of input 'Hamiltonian' H
       in QSpace represenation. The eigendecomposition is split into
       kept (K) and discarded (T) state space, as governed by the
       truncation parameters Nkeep and Etrunc (for actual Hamiltonians)
       or Rtrunc (for density matrices R:=H).
 
   Options
 
      'Nkeep',..  max. number of states to keep (default: -1 = keep all)
      'Etrunc',.. keep states according to energy truncation criteria
                  with E<=E0+Etrunc (default: <=0 = ignore; note that
                  Etrunc is interpreted relative to lowest eigenvalue
                  E0 = min[eig(H)]).
      'Rtrunc',.. keep states according to truncation criteria of
                  reduced density matrix with eig(R)>Rtrunc with R:=H
                  (this is alternative to Etrunc and thus ignored by
                  default; note, furthermore, that here Rtrunc is
                  interpreted absolute, i.e. as is).
      'deps',..   allow to keep (close to) degenerate spaces together
                  (default: 0, i.e. sharp truncation e.g. with Nkeep)
 
   Output
 
       E contains the full sorted eigenspectrum (1st column),
       and in the presence of non-abelian symmetries, the overall
       multiplet degeneracy as a 2nd column.
 
       The output/info structure I contains the fields
 
        .AK, .AD  for kept and discarded state space
        .EK, .ED, for kept and discarded eigen spectrum (diag)
        .DB       block dimension of input Hamiltonian H in terms
                  of multiplets (1st col) and states (last col).
        .NK       actual number of kept multiplets / states.
        .Etrunc or
        .Rtrunc   actual truncation treshold used
 
   AWb (C) Sep 2006 ; Mar 2013 ; Aug 2015
\end{verbatim}

\pagebreak[3]


\mymex{getDimQS} % l=12/49

\begin{verbatim}
   Usage: [D[,D2]] = getDimQS(A)
 
       get full dimension of given QSpace object.
 
       In case of non-abelian symmetries, usage #2 returns
       number of multipets as D, and the full state space
       dimension as 2nd argument D2.
 
       NB! if no second output argument (D2) is requested,
       this routine returns D if D==D2 (e.g. which is always
       the case of abelian symmetries) but [D;D2] if D!=D2.
 
   (C) AW : May 2006 ; Oct 2014
\end{verbatim}

\pagebreak[3]


\mymex{getIdentityQS} % l=56/8

\begin{verbatim}
   Usage 1: A=getIdentityQS(A [,i1, perm,'-0']);
 
      get plain identity operator from given QSpace.
 
    '-0'  this generates the 1J symbol, i.e. an identity operator
          with all indizes inward. This can be used to "revert
          arrows", and subsequently also to avoid dragging along
          a scalar singleton index if the total symmetry of a
          given wave function has q=0.
 
   Usage 2: A=getIdentityQS(A [,i1 [,at]], B [,i2 [,bt] itag, perm]);
 
      get tensor product space of input spaces defined by
      A and B. If i[12] is not specified, rank-2 objects
      are assumed taking (the first) two indizes.
 
      Default index order of the output space is (1,2,12*)
      with the itags of A and B inherited if present.
      An <itag> name for the combined output space (dim3) may be
      specified (conj flag if present, will be ignored).
 
      [06/01/2019] In addition, for an input A-tensor with LRs
      index order convention, its itags may be inherited to C
      by specifying  xor <bt> (see usage above) in the format
      '-A[:..]' where `..' indicates extra characters to be
      catenated at the end of the itag for the fused index.
 
      For example, consider an A-tensor with itags {K01,K02*,s02},
      and E a local identity operator without itags; then in order
      to ensure a complete local state space, one may write
      >> getIdentityQS(A,1,'-A:~',E)
      which generates a QSpace with itags {K01,s02,K02~*}.
      K01 is inherited from A via `A,1'; the remaining two itags
      are derived from the remaining two itags in A;
      therefore `A,i' must have i=1 or 2 with in LRs index order.
      conj-flags are properly adjusted as needed.
      Conversely, getIdentityQS(A,2,'-A:~',E)
      generates itags {K02,s02,K01~*}, e.g. used for R->L sweep.
 
   Finaly, the output index order may be changed by specifying
   a permutation perm to be applied on the final object.
   If itag is not specified, perm must be in numerical format,
   otherwise a (compact) string format is also accepted.
 
   Further options
 
     '-h'  display this usage and exit
     '-v'  verbose
 
   Note that mex files do not allow to return class objects,
   hence A is returned as a QSpace structure. To get a class
   object, use QSpace(getIdentityQS(...)), or equivalently,
   if getIdentity() is properly defined as a wrapper routine
   within matlab's Class/@QSpace (see MPS Pack), having a QSpace
   input A, this may be shortended to getIdentity(A,...).
 
   (C) AW : Apr10 ; Oct14 ; May17
\end{verbatim}

\pagebreak[3]


\mymex{getQDimQS} % l=12/11

\begin{verbatim}
   Usage: [q,dd,dc] = getQDimQS(A,dim)
 
       get dimension for every Q in QIDX for given QSpace object
       dim is dimension to check (dim=='op' can be used to
       consider both dimensions of an operator simultaneously)
 
   Data returned
 
       q    symmetry labels present (unique; n x nsymlabels)
       dd   corresponding dimension of reduced data (n)
       dc   corresponding dimension of Clebsch-Gordan (n x nsym)
 
   AW (C) Jun 2007 ; Wb,Oct12,10
\end{verbatim}

\pagebreak[3]


\mymex{getRC} % l=5/29

\begin{verbatim}
  Usage: [...]=getRC(...)
 
      C=getRC(cgr)  get RCStore data as specified by input
      C=getRC(<sym>,'--info');
 
   Wb,Oct02,15 ; Wb,Nov06,16
\end{verbatim}

\pagebreak[3]


\mymex{getSmoothSpec} % l=36/40

\begin{verbatim}
   Usage: [omega, As [,Info]] = getSmoothSpec(om, Araw, 'eps',.. [, opts])
 
     om         omega (bins) for Araw
     Araw       raw discrete spectral data; it may contain several set
                in separate columns, e.g. like data for different spins
 
   Options / Flags
 
    'om',..     new discretized set of omega values (also returned as omega)
    'sigma',..  sigma used in log-Gauss broadening (0.6)
    'eps',...   smoothly replaces log-Gauss with regular Gaussian;
                in absolute units since this routine does not know
                anything about temperature(!). For the same reason,
                eps must be specified (typically may choose eps=T).
 
    'sigma2',.. Gaussian width sigma2*eps for w<<eps (0.5)
    'alpha',..  shift in log-Gauss broadening
                sigma/4: norm (=height) preserving (default)
                sigma/2: width preserving
 
    'func'      input data is not discrete but already function of omega
    'nlog',..   number of bins per decade of omega scale
                for logarithmic discretization   (128 = nlog_fdmNRG/2!)
    'emin',..   minimum energy for omega binning (1E-8)
    'emax',..   maximum energy for omega binning (10.)
 
    'disp',..   verbose for disp>0 (1)
 
                Note that the first (last) bin contains all data below
                (above) emin (emax).
   Output
 
      om        omega corresponding to the energy binning [-emax,+emax]
      As        smoothened spectral function
 
   see also fdmNRG.
   AWb (C) May 2006
\end{verbatim}

\pagebreak[3]


\mymex{getSmoothTDM} % l=44/23

\begin{verbatim}
   Usage: [ At [,Info]] = getSmoothTDM(tt, om, Adisc [, opts])
          [ At [,Info]] = getSmoothTDM(tt, ARAW, [, opts])
 
     tt         discretized set of time dependent data
     om         omega binning for Araw
     Adisc      raw TDM data with dimensions [ nom, nop [,nnrg]]
                with nom=length(om), nop=number of operators and,
                optionally, nnrg the number of NRG iterations
     ARAW       alternativly, instead of om and Adisc, the RAW data
                returned by tdmNRG() in partial mode can be specified
                as input
 
   Options / Flags
 
    'disc'      discrete data (i.e. does not require d(om) weight; default)
    'func'      functional data (i.e. DOES require d(om) weight)
                default: disc for rank(A)==3 and func for rank(A)==2
    '-q'        quiet mode (reduce number of log output)
 
    'alpha',..  exponential decay of energies exp(-alpha dE) [0.]
                this applies alpha in time domain (F. Anders 2005)
    'sigma',..  applies broadening in frequency domain (log.gauss)
                NB! either alpha or sigma can be specified but not both
                default: sigma=0.1 (with alpha ignored)
    'Lambda',.  Lambda from Wilson disretization (required with alpha)
    'eps',...   smoothly replaces log-Gauss with regular Gaussian;
                in absolute units since this routine does not know
                anything about temperature(!). For the same reason,
                eps must be specified (typically may choose eps=T).
 
    'nlog',..   number of bins per decade of omega scale
                for logarithmic discretization   (100)
    'emin',..   minimum energy for omega binning (1E-8)
    'emax',..   maximum energy for omega binning (10.)
 
                Note that the first (last) bin contains all data below
                (above) emin (emax).
   Output
 
      At        input spectral information transformed into time
                dependent data.
      Info      structure with additional information
 
   see also tdmNRG.
   AW (C) Jan 2007
\end{verbatim}

\pagebreak[3]


\mymex{getSymStates} % l=37/14

\begin{verbatim}
   Usage #1: [U,I]=getSymStates(Sp,Sz);
 
     get symmetry eigenstates for symmetry operations given by
     rising operators Sp and (diagonal) z-operators Sz.
     The order of the Sz operators determines the order of the
     z-labels to be determined.
 
   Usage #2: [U,I]=getSymStates(SOP)
 
     Moreover, Sp and Sz can be grouped according to each symmetry
     by using a single vector structure Is with the mandatory
     fields .Sp and .Sz. In this case, an empty Sp is interpreted
     as Abelian quantum number. Eventually Sp and Sz are combined
     into a single sequence as in usage #1.
 
     Further optional fields
       .type   symmetry type (in QSpace convention)
       .qfac   applied to all z- and q-labels (scalar or vector)
       .jmap   subsequently to qfac, applies linear transform
               of q-labels (square matrix)
 
   Output: symmetry eigenstates written as full matrix in U,
   and info structure I with the fields
 
     basically reiterating the input
       .R0    input .Sp and .Sz data
       .type  symmetry type
       .JMap  J-transformation used (identity if not specified)
       .dz    number of z-operators for each symmetry
 
     actual output data
       .dd    dimension of irreduible multiplets generated
       .RR    irreducible representations with qfac and jmap applied
       .QZ    interleaved q- and z-labels.
       .d2    dimensions of setup in QZ (non-abelian number have
              the same number appearing twice)
 
   (C) Wb,Jul05,10
\end{verbatim}

\pagebreak[3]


\mymex{isIdentityCG} % l=5/57

\begin{verbatim}
   Usage: i = isIdentityCG(A)
 
       check whether QSpace A has indentity Clebsch-Gordan spaces.
       Note that for full abelian spaces, this is always true.
 
   Wb,Jan12,12
\end{verbatim}

\pagebreak[3]


\mymex{isIdentityQS} % l=6/9

\begin{verbatim}
   Usage: i = isIdentityQS(A [,eps])
 
       check whether QSpace A represents a block-diagonal
       identity matrix. This also checks the corresponding
       Clebsch-Gordon spaces if any (default: eps=1E-12).
 
   (C) AW : May 2006 ; Oct 2010 ; Oct 2014
\end{verbatim}

\pagebreak[3]


\mymex{makeUniqueQS} % l=9/21

\begin{verbatim}
   Usage: A = makeUniqueQS(A [,opts])
 
       MEX wrapper routine for QSpace::makeUnique()
 
   Options
 
       -[qQ]    quite mode
       -[vV]    verbose mode (wrt. outer multiplicity)
 
   AWb (C) Nov 2007 ; 2012
\end{verbatim}

\pagebreak[3]


\mymex{mpsIsHConj} % l=5/36

\begin{verbatim}
   Usage: [i,istr]=mpsIsHConj(A [,eps]);
 
   check whether QSpace is hermitian within threshold eps (1E-12).
   A must be even-rank object as hyperindex is allowed.
 
  Wb,Aug17,06
\end{verbatim}

\pagebreak[3]


\mymex{normQS} % l=7/47

\begin{verbatim}
   Usage: nrm = normQS(A)
 
       calculates the (Frobenius) norm of given QSpace,
       i.e. |A|^2 = norm(A)^2 = trace(A'*A), and appropriately
       generalized to arbitray rank. In the presence of CGCs, this
       implies |A|^2 = sum_i (|data{i}|^2 * prod_j(|cgs(i,j)|^2)).
 
    (C) AW : Jun 2010 ; Oct 2014
\end{verbatim}

\pagebreak[3]


\mymex{orthoQS} % l=71/1

\begin{verbatim}
   Usage #1:
   [A1,A2,[,info]] = orthoQS(PSI,idx, [,OPTS]);
 
      Orthonormalize QSpace PSI w.r.t. given SINGLE index idx,
      which is split off into tensor A1, connected via a trailing
      new intermediate index with A2, i.e. PSI=A1*A2, where
      A2 keeps its index order (i.e. index idx stays in place).
 
      Note that the new intermediate index will typically be
      truncated w.r.t. to singular values (see stol below).
 
      Convention on orthonormalization direction (cf. lrdir below,
      which is NOT specified here): A1 is ALWAYS an isometry/unitary
      that has one in- and one out-index. This allows the new
      intermediate index to carry the same itag as idx in PSI.
 
   Usage #2 (abelian symmetries only)
   [A1,A2,[,info]] = orthoQS(PSI,idx, lrdir [,OPTS]);
 
      Similar to usage #1 above, except that for abelian symmetries,
      idx may contain any number of indices consistent with PSI
      (in constrast, for non-abelian symmetries, idx can only
      specify a single index (or equivalently r-1 indices with
      r the rank of PSI.
 
   Usage #3 (abelian symmetries only)
   [A1,A2,[,info]] = orthoQS(A1,A2,ic1,ic2, lrdir [,OPTS]);
 
      Orthonormalize QSpace A1 towards index ic1 using SVD,
      where the resulting tensor, split off from A1, is right
      away merged with A2 at index ic2 (having lrdir='>>';
      vice versa for lrdir='<<').
 
      This internally connects to usage #2, in that the input tensors
      A1 and A2 are right away contracted; e.g. for lrdir='>>':
      PSI = contractQS(A1,ic1,A2,ic2);
 
   INPUT
 
      PSI      or PSI=A1*A2, where
      A[12]    are two nearest neighbor tensor
      ic[12]   that share the common index specified by ic1 and ic2.
 
      lrdir    direction of orthonormalization; this accepts the
               following equivalent values:
               '>>', 'LR', or +1 for "left-to-right"
               '<<', 'RL', or -1 for "right-to-left"
   OPTIONS
 
     'Nkeep',.. max. number of states/multiplets to keep (-1: all)
     'Nkmin',.. min. number of states/multiplets to keep on index
                connecting A1 to A2 (default: 0, i.e. no lower bound)
     'stol',... absolute tolerance on SVD i.e. on sqrt(rho)
                default: 1E-8 if Nkeep is not set, 0 otherwise.
     'itag',... itag on newly generated index
     '-v'       print info
 
   OUTPUT
 
      A1[2]   orthonormalized A1[2]
      info      info structure containing (amongst other fields)
        .svd    vectorized sorted complete singular value data;
                multiplicity (i.e. multiplet dimension) is specified
                in svd(:,2) in the presence of non-abelian symmetries.
        .svd2tr truncated weight (singular values squared).
        .sfac   factor to reestablish the orginal norm(PSI).
        .flag   !=0 indicating potential trouble (see .info field).
        .DD     (multiplet) dimension of original PSI.
        .D1     (multiplet) dimension of A1.
 
   See also svdQS.
   AW (C) Jun 2006; Jun 2009; Aug 2015
\end{verbatim}

\pagebreak[3]


\mymex{permuteQS} % l=18/19

\begin{verbatim}
   Usage: A = permuteQS(A, P [,'conj'])
 
       permute input QSpace using given permutation P.
 
       Optional trailing 'conj' also applies (complex) conjugation
       (note that this also affects real QSpaces in that qdir and
       itags are altered!).
 
       For convenience, P [,'conj'] may also be represented as
       single string, e.g. [2 1],'conj' is equivalent to '2,1;*'
       or '21*' where the convention on string notation
       follows that of contraction indices [ctrIdx].
 
       NB! [06/02/2019] the provided permutation can be shorter
       than the rank of the QSpace; in this case it only affects the
       leading range of indices, i.e., acts like an identity
       on the remainder of indices.
 
   AW (C) Aug 2006 ; May 2010 ; Oct 2014
\end{verbatim}

\pagebreak[3]


\mymex{plusQS} % l=9/43

\begin{verbatim}
   Usage: C = plusQS(A, B [, bfac, opts])
 
       adds one QSpace to another as in C = A+bfac*B ([bfac=1])
       block-diagonal identity matrix.
 
   Options
 
      '-v'  verbose mode (e.g. on remaining outer multiplicity)
 
   (C) AW : Jun 2006 ; Oct 2014
\end{verbatim}

\pagebreak[3]


\mymex{setupRCStore} % l=19/58

\begin{verbatim}
   Usage #1: setupRCStore(sym [,opts])
 
       Setup CGC container for given symmetry,
       with the data stored in directory specified by the
       environmental variable $RC_STORE.
 
   Usage #2: setupRCStore(sym,q1,q2 [,opts])
 
       explicitely generate tensor product decomposition q1*q2
 
   Options
 
      'npass',... number of passes through gCG data (3).
      'dmax',...  specify max. average multiplet dimension sqrt(d1*d2)
                  in tensor product decomposition (10*qlen).
      '-a'        generate all, yet compatible with dmax (by default,
                  multiplets with 2-digit qlabels are no longer used
                  in decomposition)
 
   (C) Wb,Aug07,14 ; Wb,Apr01,15
\end{verbatim}

\pagebreak[3]


\mymex{skipZerosQS} % l=11/24

\begin{verbatim}
   Usage: [A,n] = skipZerosQS(A [,eps,opts]);
 
      skip data blocks that have all-zero entries to within eps
      (default eps=1E-14). The total number of skipped blocks
      is returned as 2nd argument.
 
      For scalar (i.e. rank-2 block-diagonal) operators zero-
      blocks are kept, by default (e.g. relevant for Hamiltonians
      where states with energy zero are relevant states.
      This can be disable by using the option '--all'.
 
   Wb,May14,10
\end{verbatim}

\pagebreak[3]


\mymex{svdQS} % l=58/41

\begin{verbatim}
   Usage #1:
   [U,S,Vd,[,info]] = svdQS(PSI,idx [,OPTS]);
 
      Compute SVD decomposition of QSpace object PSI, where the
      indices idx are split off into Vd, and the remaining indices
      into U, which are '>>' and '<<'  orthonormalized, respectively
      (i.e. are unitaries or isometries generated from the SVD.
      Typically, the indices in PSI are 'all-in'.
 
      Note that the 3rd argument returns Vd:=V' and not V,
      i.e. already includes the dagger such that PSI = U*S*Vd
      [for idx=2; this is also LAPACK convention; for idx=1
      the transpose is returned for U, i.e., PSI = [(U^t)*S*Vd]^t,
      since, by convention, U resembles the index order of PSI].
 
      For non-abelian symmetries, idx is required to be a single
      index (or conversely, all indices but one). If PSI is grouped
      into two sets of indices each of which has more than one index,
      then indices need to be fused prior to calling svdQS.
      Also, in this case of a single index specified, U always
      resembles PSI, i.e. has the same rank and index order.
 
   Usage #2 (deprecated):
   [U,S,V',[,info]] = svdQS(A1, A2, ic1, ic2 [,OPTS]);
 
      Compute SVD decomposition of QSpace object PSI=A1*A2,
      where '*' stands for the contraction (A1,ic1,A2,ic2).
      The split occurs between the remaining indices in A1
      and the remaining indices in A2. This indices are typically
      'all-in', while the index/indices ic1 to be contracted
      with ic2 can have any initial orientation. There will
      always be a single new intermediate connecting index
      between A1 and A2.
 
   Input
 
      PSI      or PSI=A1*A2, where
      A[12]    are two nearest neighbor tensor
      ic[12]   that share the common index specified by ic1 and ic2.
 
   OPTIONS
 
     'Nkeep',.. max. number of states/multiplets to keep (-1: all)
     'Nkmin',.. min. number of states/multiplets to keep on index
                connecting A1 to A2 (default: 0, i.e. no lower bound)
     'stol',... absolute tolerance on SVD i.e. on sqrt(rho)
                default: 1E-8 if Nkeep is not set, 0 otherwise.
     'itag',... itag on newly generated index
     'disp'     display info to stdout
 
   OUTPUT
 
     A1(2)   updated and properly orthonormalized A1(2)
     info      info structure containing
      .SVD     singular value data blocked within QS symmetry space
      .svd     vectorized singular value data
 
   See also orthoQS.
   AW (C) Sep 2008 ; Aug 2015
\end{verbatim}

\pagebreak[3]


\mymex{traceQS} % l=8/46

\begin{verbatim}
   Usage: x = traceQS(A [,i1,i2])
 
       trace out index set i1 with i2 of given QSpace object;
       for even-rank A, missing index sets is interpreted as
       full trace of indices i1=1:r/2 with iw=r/2+1:r.
       if all indices of A are traced out, the returned x is
       a number (QSpace otherwise).
 
   AW (C) 2006 ; 2018
\end{verbatim}

\mysec[MEX routines (util)]{MEX routines (util)}


\pagebreak[3]


\mymex{filestat} % l=16/1

\begin{verbatim}
  Usage: i=filestat(f,tag [,'-s'])
 
     get system file information
 
     If '-s' is specified, a time string is returned;
     otherwise (by default) the returned date is in Matlab datenum
     format, e.g. can be used with datestr(t,*).
 
  Available tags
 
    'm'     modification time stamp
    'c'     creation time stamp
    'a'     last access time stamp
 
  All these time stamps return split-second resolution.
 
  Wb,Jan28,16
\end{verbatim}

\pagebreak[3]


\mymex{getchar} % l=4/23

\begin{verbatim}
   Usage: getchar()
 
      C++ program that reads single character from keyboard
 
   Wb,Jan09,09
\end{verbatim}

\pagebreak[3]


\mymex{kramers} % l=25/33

\begin{verbatim}
   Usage: [yr,I]=kramers(xi,yi [,OPTS])
 
     calculate Kramers-Kronig transform from discrete data
     using linear interpolation scheme.
 
       xi         x-data
       yi         y-data
       yr         kramers kronig transformed data (default: imag->real)
 
   Options
 
      '-disc'     yi data is already discrete
      'eps',..    for continuous data: regularization of principal
                  value data for xi==xj (xi += eps*(xi-x[i-1]; 1E-10)
    If isdisc
 
      'faceta',.. width of Lorentzian delta function relative to
                  data point spacing dx (1.)
   else
 
      'g1flag'    include tails to infinity by assuming that
                  input yi decays like 1/x for large x
      'g2flag'    same as g1flag but assuming 1/x^2 decay
      'r2i'       do real->imag transform instead (=overall minus sign)
 
   Wb,Dec12,06
\end{verbatim}

\pagebreak[3]


\mymex{matchIndex} % l=24/5

\begin{verbatim}
   Usage: [Ia,Ib,I] = matchIndex(QA, QB [,OPTS])
 
      QA or QB are matrizes whos rows are matched; consequently, 
      their number of columns (row dimension) must be equal.
 
   Options
 
       '-s'    sort indizes with respect to A
       '-q'    quiet mode, i.e. do not issue warning when NaN's
               are encountered.
       'eps',..specify epsilon within which to accept match (0)
       '-f'    use float instead of double (gets rid of numerical
               noise in case of double precision input data)
 
   Output
       Ia      index into QA (QA.Q)
       Ib      index into QB (QB.Q)
 
       I       an info structure with elements
       .n      number of matches [=length(Ia)]
       .dim    number of input records [ size(QA,1), size(QB,1) ]
       .ix1    entries in QA not matched in QB
       .ix2    entries in QB not matched in QA
 
   AW (C) 2006 ; Feb10,11
\end{verbatim}

\pagebreak[3]


\mymex{mexworld} % l=1/35

\begin{verbatim}
   This is a help on the mex-file mexworld
   Wb,Jan12,19
\end{verbatim}

\pagebreak[3]


\mymex{mpfr2dec} % l=10/42

\begin{verbatim}
   Usage: [data vector]=mpfr2dec([char matrix] [,base]);
 
      Convert mpfr data strings to decimals.
      Since matlab is col-major, the data string for each
      value is expected to represent an entire column.
 
   Example
 
      mpfr2dec(['0':'9', 'A':'Z', 'a':'z'],62)
 
   Wb,Jul01,14
\end{verbatim}

\pagebreak[3]


\mymex{printfc} % l=9/58

\begin{verbatim}
   Usage: printfc(fmt)
 
      C++ program that emulates C printf
      ie. does not terminated automatically with a newline and
      allows for escape sequences (in contrast to fprintf).
 
      NB! If more than one onput argument is specified, all
      arguments are forwarded to matlab's sprintf.
 
   AWb (C) Jan 2006
\end{verbatim}

\pagebreak[3]


\mymex{sparse2wb} % l=4/14

\begin{verbatim}
   function S=sparse2wb(S)
 
      convert matlab sparse array to wbsparse format.
 
   AW (C) 2011
\end{verbatim}

\pagebreak[3]


\mymex{uniquerows} % l=11/24

\begin{verbatim}
   Usage: [A [,I,D]] = uniquerows(Q [,opts])
 
      MEX file that makes rows of array unique. the second return
      argument is a cell array of index sets of rows with the same
      data. D specifies the corresponding multiplicity of the
      unique rows.
 
   Options
 
     '-1'  return only first index in each group in I
 
   Wb (C) Sep 2006
\end{verbatim}

\pagebreak[3]


\mymex{wbeig} % l=5/41

\begin{verbatim}
   Usage: [E,U] = wbeig(H);
 
       eigenvalue decomposition based on BLAS dsyev
       for (genearlized) symmetric rank-2 object.
 
   AW (C) 2009
\end{verbatim}

\pagebreak[3]


\mymex{wbhash} % l=34/52

\begin{verbatim}
   Usage: [kk,I] = wbhash(dd,N [,opts]);
 
      Compute keys kk using given hash for each row of the
      input array using given hash specified by the hash
      function (uses offset and multiplier) together with
      the final dimension N of the hash key (hash mod N).
      The hash function is given by
 
          unsigned long h=offset;
          for (i=0; i<len; ++i) {
             h = mult*h + s[i];
          }
 
      Here s represents a row in dd of length len.
      Default values for the hash function are
 
          offset = 5381 (prime)
          mult = 33 = 2^5 + 1;
 
      which correspond to the hash function 'djb2'
      as suggested by Daniel B. Bernstein Having mult=33,
      the hash value is evaluated through h = (h<<5)+h+s[i].
 
   Options
 
    For algorithm based on multiplier
      'mult',..   multiplier (33)
      'offset',.. offset (5381)
 
    Alternatively the following flags use default string hashes
      '-stl'      uses hash<std::string> string_hash(char(dd)).
      '-wbb'      adapted versions of algorithms in boost/hash.cpp
      '-ump'      test mode for unordered_map
 
   Wb,Jul01,14
\end{verbatim}

\pagebreak[3]


\mymex{wbhist} % l=31/33

\begin{verbatim}
   Usage: ybin=wbhist([strid,] x,y [,xbin]);
 
      mex routine to collect binned data. If y is not specified,
      this generates a simple histogram counting occurances.
 
      NB! The binning occurs with respect to the mid-points of
      the xbin data, where the first bin contains everything for
      x < mean(xbin(1:2)), and the last bin everything for
      x >= mean(xbin(end-1:end)).
 
      NB! data is stored globally within mex file, i.e. is
      cleared when using 'clear functions'. However, this allows
      subsequent calls without last argument xbin to add to
      existing binned data. The last argument must be specified
      for the first call given some strid, and whenever the
      histogram data shall be initialized.
 
      default hist set has strid 'default'.
 
   Further tags:
 
     'stat'         current status of wbhist (global)
     'clear'        clear entire wbhist (global)
     'id,'remove'   remove wbhist data set associated with 'id'
 
   Examples
 
      wbhist([],[],xbin);
      wbhist(x, []);   plain count (assumes y=ones(size(x))).
      [a,b,c]=wbhist(['default']);
 
   Wb,Feb12,11
\end{verbatim}

\pagebreak[3]


\mymex{wbrat} % l=27/11

\begin{verbatim}
   Usage: [q,I] = wbrat(q [,opts]);
 
      get rational approximation of given number
      based on continued fraction approach.
 
   Options
 
      '-h'      display this usage
      '-v'      verbose mode
      '-q'      quiet mode
      '-r'      assume also square root data (using rat() on q^2)
      '-a'      also return continued fraction sequence in I
      'nmax',.. maximimum number of iterations (nmax)
      'eps1',.. tolerance on continued fraction (1E-6)
      'eps2',.. rel. overall tolerance wrt. to in put number q (1E-12)
 
   Output
 
       q        same as input, but 'fixed' to rational approximation
       I        info structure
                .P  integer enumerator
                .Q  integer denominator
                .relerr relative error of rational representation
                .n  total number of continued fraction iterations used
                .param  (input) parameters nmax, eps1, eps2 used
 
   See also MatLab's rat.m
   AWb (C) July 2010
\end{verbatim}

\pagebreak[3]


\mymex{wbsvd} % l=9/44

\begin{verbatim}
   Usage: [U,S,V] = wbsvd(H [,I]);
 
       singular value decomposition based on BLAS dgesvd
       for genearlized rank-r object.
       index set I specifies which dimensions are combined
       into first fused dimension. For rank-2 objects
       and I omitted, this corresponds to the regular SVD
       on matrizes.
 
   AW (C) Sep 2009
\end{verbatim}

\newpage


\mysec[QSpace member functions]{Class/@QSpace member functions}


\pagebreak[3]


\myfun{QSpace/Casimir.m} % l=1/0

\begin{verbatim}
  function Casimir(S)
  Wb,Oct05,15
\end{verbatim}

\pagebreak[3]


\myfun{QSpace/QSpace.m} % l=36/7

\begin{verbatim}
  Function [A,varargout] = QSpace(...)
 
     Construct QSpace class object.
 
     This serves to wrap QSpace structures into QSpace objects,
     and also can be used to generate several specific QSpace objects.
     For historical reasons, this also includes low-level QSpace
     contructors for the case of only abelian quantum numbers.
 
  Usage:
 
  #1: A = QSpace(Q,data)
      A = QSpace(A)   conversion e.g. from struct
      A = QSpace(A,'check')  also performs full check on QSpace structure
 
  #2: A = QSpace(q1, data1, q2, data2, ...) % see initQSpace
 
  #3: A = initQSpace(Q1,M,DB)
          square operator M with row/col Q1 block size given by DB
 
  #4: A = initQSpace(Q1,Q2,M,'map')
          matrix M with row/col QIDX given by Q1/Q2
          eg. used for projectors
 
  #5: A = QSpace(Q,M,'operator');
  #6: A = QSpace(A,'unity');
 
  #7: A0=initQSpace(QL,Qs1,Qs2,...); setup identity tensor
          for product spaces (QL,Qs1,...)
   
          it must hold: size(QL,2) = size(Qs1,2) = size(Qs2,2) = ... = QDIM.
          Qs and QL span the full Hilbert space, hence certain quantum
          lables may appear several times.
 
  Wb,May12,06 ; Wb,Apr16,16

    Documentation for QSpace
\end{verbatim}

\pagebreak[3]


\myfun{QSpace/SEntropy.m} % l=15/49

\begin{verbatim}
  function [se,rr,I]=SEntropy(Rho [,eps,opts])
 
    Calculate von Neumann entropy of the mixed (reduced)
    density matrix R.
 
    R may also contain non-abelian symmetries in which case
    the output weights rr are already weighted by their 
    corresponding multiplet degeneracies such that sum(rr)=1.
 
  Options
 
     eps      check trace(Rho)=1 within numerical accuracy of eps (1E-12)
    '-xm'     expand multiplet degeneracy in returned rr
    'beta',.. assume R is a Hamiltonian => apply beta=1/T and exponentiate
 
  Wb,Jan24,11
\end{verbatim}

\pagebreak[3]


\myfun{QSpace/SetCC.m} % l=5/11

\begin{verbatim}
  Function [A,t]=SetCC(A [,idx])
 
     Set "dagger" for specified index-labels (all if idx is not specified).
     If any A.data{k} is complex, A.data will also be complex conjugated.
 
  Wb,Aug04,12
\end{verbatim}

\pagebreak[3]


\myfun{QSpace/SetOL.m} % l=11/22

\begin{verbatim}
  function A=SetOL(A,X,k [,otag])   % usage #1
  function A=SetOL(A,itag[,otag])   % usage #2
 
     Set operator labels of the type { itag, itag* [,'op*'] }.
     default operator tag (otag) for rank-3 tensors is 'op'.
 
  Usage #1
 
     set index labels for operator A based on the label found
     in tensor X at dimension k.
 
  Wb,Aug04,12
\end{verbatim}

\pagebreak[3]


\myfun{QSpace/acomm.m} % l=5/39

\begin{verbatim}
  function C=acomm(A [,B,cflag])
 
     cflag if set, must equal 'conj'
     and is interpreted as conjA, i.e. C = {A',B}
 
  Wb,Feb06,16
\end{verbatim}

\pagebreak[3]


\myfun{QSpace/appendScalarSymmetry.m} % l=20/50

\begin{verbatim}
  function A=appendScalarSymmetry(A,qtype [,opts])
 
     Append scalar symmetry (e.g. new symmetry label,
     where all existing states get a newly added qlabel q=0
     with the zeros repeated according to the rank of the
     added symmetry.
 
  Options
 
    'pos',..  actually insert symmetry at particular position
              in symmetry order; e.g., 'pos',1 will prepend the symmetry.
    'q',..    specify particular symmetry labels other than (default q=0).
              if non-zero, this assumes A to be of operator type in that
              specified q is applied to the first two indices only *)
 
  *) rather than specifying non-zero q, consider starting with zeros
     first, followed by mapping of operators into tensor product
     of state spaces of operators
 
  See also: 
  Wb,Jul12,11 ; Wb,Dec08,14
\end{verbatim}

\pagebreak[3]


\myfun{QSpace/appendSingletons.m} % l=23/17

\begin{verbatim}
  function A=appendSingletons(A,'new_itag[*]')
  function A=appendSingletons(A,qdir)
  function A=appendSingletons(A,r_new,qdir)
 
     Add singleton dimension(s) to reach rank r_new.
     qdir sets direction of new singleton dimensions (numeric; e.g. +/-1)
 
     In usage #1 assigns the specified itag to the newly added singleton
     dimension; this also permits to simultaneously specify qdir via a
     trailing conjugation flag '*'.
 
     If numel(qdir)>1, it is assumed it specifies qdir for all of A
     where the values in qdir are ignored for the ones already existing
     in A [this is useful if qdir was obtained via getqdir(Aother)].
 
     NB! use makeIrop() to append irop dimension instead of this,
     since there may be non-zero symmetry labels for all-abelians ops.
 
  Examples
 
     A4=appendSingletons(E,'  +-')
 
  See also makeIrop, fixScalarOp.
  Wb,Sep17,18
\end{verbatim}

\pagebreak[3]


\myfun{QSpace/applyK1K2.m} % l=7/46

\begin{verbatim}
  function [S1,S2]=applyK1K2(S,H)
 
     auxilliary routine that applies energy kernels
     K1 (principal value) and K2 (remainder for diagonal values).
     e.g. required for static magnetic susceptibility
     (K2 accounts for correction due to degeneracies).
 
  Wb,Jul05,12
\end{verbatim}

\pagebreak[3]


\myfun{QSpace/applyZ.m} % l=15/0

\begin{verbatim}
  function A=applyZ(A,idx [,iq,opts])
 
     apply fermionic signs to given index / indices (idx)
     with fermionic charge (or parity Z2) assumed at position
     iq in the Q-labels (iq may be skipped, if only one
     abelian symmetry label is present).
 
  Options
 
    'qfac',.. : factor to apply on charge labels, to obtain
     actual charge (must be integer; enforced); e.g. note that
     single-flavor spinless fermions have an artificial factor
     of 2 on their charge labels to permit simple symmetry around
     half-filling; hence in this case: applyZ(...,'qfac',1/2);
 
  Wb,Sep17,18
\end{verbatim}

\pagebreak[3]


\myfun{QSpace/blkdiag.m} % l=1/21

\begin{verbatim}
  function M=blkdiag(A,varargin)
  Wb,Aug10,16
\end{verbatim}

\pagebreak[3]


\myfun{QSpace/cat.m} % l=7/28

\begin{verbatim}
  function X=cat(A1,A2,...,k)
 
     Concatenate data sectors of given QSpaces along dimension k
     (QSpace vectors are also included in the concatenation procedure).
     NB! k may be a cell, i.e. {k,qdir} to specify qdir of newly added
     dimensions.
 
  Wb,Sep17,18
\end{verbatim}

\pagebreak[3]


\myfun{QSpace/cgr\_size.m} % l=1/41

\begin{verbatim}
  function s=cgr_size(A,i,j)
  Wb,Jan08,12
\end{verbatim}

\pagebreak[3]


\myfun{QSpace/checkOM.m} % l=4/48

\begin{verbatim}
  function m=checkOM(A [,opts])
 
     Check for outer multiplicity in given QSpace
 
  Wb,Feb15,14
\end{verbatim}

\pagebreak[3]


\myfun{QSpace/checkSameSpace.m} % l=10/58

\begin{verbatim}
  function i=checkSameSpace(A,ia,B,ib)
 
     check whether QSpaces A and B share the same vector spaces
     (simple check in terms of symmetry blocks and respective
     multiplet [and Clebsch-Gordan dimensions, if applicable]).
 
     This routine uses getQDimQS(A,ia) [and similar for B]
     hence ia accepts the same index specification as getQDimQS
     that is, a specific dimension or 'op'(erator).
 
  Wb,Jan12,12
\end{verbatim}

\pagebreak[3]


\myfun{QSpace/checkdim.m} % l=1/15

\begin{verbatim}
QSpace/checkdim is a function.
    [ok, ii] = checkdim(A, pp)
\end{verbatim}

\pagebreak[3]


\myfun{QSpace/chopd.m} % l=6/22

\begin{verbatim}
  function A=chopd(A [,aref])
 
     apply chopd() to data
     => sets small matrix elements to zero
     => removes imaginary noise if any
 
  Wb,Oct17,11
\end{verbatim}

\pagebreak[3]


\myfun{QSpace/comm.m} % l=13/34

\begin{verbatim}
  function C=comm(A [,B,'conj','+'])
 
     Compute commutator (or anticommutator, if '+' is specified).
     if B is ommitted, this assumes B=A && 'conj'.
     'conj' if set, is interpreted as conjA, i.e. C = {A',B}
     For anticommutator, use comm(A [, B, cflag], '+');
 
     The operators and B must be rank-2 or rank-3, where the latter
     case assumes operator index order (ss',o), i.e. the irop (=spinor)
     index comes last.
 
  Ex. comm(F,'+') computes {F',F}
 
  Wb,Feb06,16
\end{verbatim}

\pagebreak[3]


\myfun{QSpace/compact.m} % l=1/53

\begin{verbatim}
  merge quantum numbers of tensor operator
  Wb,Jul07,09
\end{verbatim}

\pagebreak[3]


\myfun{QSpace/conj.m} % l=5/1

\begin{verbatim}
  Function: A=conj(A)
 
     get complex conjugate of operator in QSpace format;
     overloading the conj routine
 
  Wb,Jun02,08
\end{verbatim}

\pagebreak[3]


\myfun{QSpace/contract.m} % l=5/12

\begin{verbatim}
  function A=contract([args])
 
     same as contractQS
     but returns QSpace object (rather than plain structure).
 
  Wb,Sep18,06 ; Wb,Mar31,16
\end{verbatim}

\pagebreak[3]


\myfun{QSpace/ctranspose.m} % l=2/23

\begin{verbatim}
  overloading the hermitian conjugate operator '
  accept hyper index
  Wb,Nov30,09
\end{verbatim}

\pagebreak[3]


\myfun{QSpace/datasize.m} % l=6/31

\begin{verbatim}
  Function [D,DD]=datasize(A [,'-s'])
 
     get full listing of size(data{}),
     except if option '-s' is specified which gives summary 
     in terms of total size in bytes of A.data
 
  Wb,Apr24,08
\end{verbatim}

\pagebreak[3]


\myfun{QSpace/decomposeOp.m} % l=6/43

\begin{verbatim}
  function [cc,err,I]=decomposeOp(H,{O1,O2,...}[,R][,opts])
 
     decompose operator H into operator (sets) specified in X={O1, O2, ...};
     if R is specified, operator overerlap is calculated
     w.r.t. density matrix R.
 
  Wb,Apr16,13
\end{verbatim}

\pagebreak[3]


\myfun{QSpace/deconj.m} % l=5/55

\begin{verbatim}
  Function [A,t]=deconj(A [,idx])
 
     Unset "bra" (complex conjugate) label for specified dimensions
     (all if idx is not specified).
 
  Wb,Apr17,13
\end{verbatim}

\pagebreak[3]


\myfun{QSpace/diag.m} % l=8/7

\begin{verbatim}
  Function: A=diag(A [,OPTS])
 
     extract diagonal of operator in QSpace format
     alternatively, if only diagonal is stored,
     diagonals become full again.
 
  Options
 
  Wb,Sep08,06
\end{verbatim}

\pagebreak[3]


\myfun{QSpace/dim.m} % l=9/21

\begin{verbatim}
  function s=dim(A [,opts])
 
  Options
 
    '-a'  all (simply return output from getDimQS())
    '-f'  full dimension
          (in case of CGC spaces, corresponds to state space dimension)
    '-op' IROP dimension (i.e. checking for rank-3 QSpace)
 
  Wb,Jul31,12
\end{verbatim}

\pagebreak[3]


\myfun{QSpace/disp\_cgr.m} % l=2/36

\begin{verbatim}
  function disp_cgr(A)
  see also plot_cgr()
  Wb,Jan13,12
\end{verbatim}

\pagebreak[3]


\myfun{QSpace/disp\_qdim.m} % l=2/44

\begin{verbatim}
  function disp_qdim(A,varargin)
    example: disp_qdim(A0,2,IS.SOP)
  Wb,Jan26,12
\end{verbatim}

\pagebreak[3]


\myfun{QSpace/display.m} % l=12/52

\begin{verbatim}
  Function display(A [,OPTS])
  Options
 
    '-a'       show all entries
    '-v'       verbose flag (shows CGS dimension separately for each symmetry)
    '-E'       sort QSpace wrt. data (energy) [assumes diagonal Hamiltonian]
    '-s'       sort QIDX
    'sperm',.. sort QIDX using given permutation
 
     By default, QSpaces with many entries (length > 14) will show
     a truncated list, including also the two largest entries.
 
  Wb,Mar01,08
\end{verbatim}

\pagebreak[3]


\myfun{QSpace/dsize.m} % l=10/11

\begin{verbatim}
  Function [D,[DD|I]]=dsize(A [,opts])
  get maximum D for every dimension of given MPS state
 
      where D is a vector of length rank(A)
 
  Options
 
    '-d'  returns detailed block sizes of single input A in DD with
          CG space considered if exists (overall size returned in D)
 
  Wb,Aug08,06
\end{verbatim}

\pagebreak[3]


\myfun{QSpace/eig.m} % l=3/27

\begin{verbatim}
  function [varargout]=eig(H,varargin)
     ee=eig(H,varargin);
     [A,E]=eig(H,varargin);
  Wb,Aug30,12
\end{verbatim}

\pagebreak[3]


\myfun{QSpace/eq.m} % l=0/36

\begin{verbatim}
  overload of == (equal) operator
\end{verbatim}

\pagebreak[3]


\myfun{QSpace/expandQ.m} % l=1/42

\begin{verbatim}
  Function qq=expandQ(A,k)
  deprecated: see getQRange.m
\end{verbatim}

\pagebreak[3]


\myfun{QSpace/finditag.m} % l=5/49

\begin{verbatim}
  function i=finditag(A,pat [,n])
 
     find itag with specit regexp pattern
     if n is specified, the result must yield exactly n matches.
 
  Wb,Jun03,19
\end{verbatim}

\pagebreak[3]


\myfun{QSpace/fixAbelian.m} % l=8/1

\begin{verbatim}
  function fixAbelian(A,B,...)
  Options
 
     -xop3    reduce "scalar" rank-3 operators to regular rank-2.
 
     -addCGR  add (trivial) CGR space to all-abelian operators.
     -op3     together with -addCGR: make operator of rank 3 if rank 2.
 
  Wb,Feb20,13
\end{verbatim}

\pagebreak[3]


\myfun{QSpace/fixScalarOp.m} % l=11/15

\begin{verbatim}
  function fixScalarOp([opts,] A,B,...)
 
    reduces "scalar" rank-3 operators to regular rank-2.
    converting (yet keeping) the Clebsch Gordan coefficient
    spaces to sparse.
 
  Options
 
    '-f'  enforce reduciblity if rank-3.
 
  See also makeIrop() to append 3rd dimension
  Wb,Aug21,13
\end{verbatim}

\pagebreak[3]


\myfun{QSpace/full.m} % l=0/32

\begin{verbatim}
  wrapper routine for mpsFullQS()
\end{verbatim}

\pagebreak[3]


\myfun{QSpace/full2.m} % l=0/38

\begin{verbatim}
  wrapper routine for mpsFull2QS()
\end{verbatim}

\pagebreak[3]


\myfun{QSpace/gdisp.m} % l=4/44

\begin{verbatim}
  function gdisp(A)
 
     graphical information regarding given QSpace A.
 
  Wb,Feb09,11
\end{verbatim}

\pagebreak[3]


\myfun{QSpace/getAtensorLoc.m} % l=15/54

\begin{verbatim}
  function A=getAtensorLoc(X [,qvac][,l])
     
     get Atensor by extracting local state space from local
     operator X together with vacuum state on virtual bond
     (default for qvac; optionally any other bond-space can
     be specified as third argument, instead).
 
     l \in {'L','R'} specifies the location of vacuum state
     using LRs index order (default: l='L').
 
     All arguments except for the last (l) are directly
     handed over to getIdentityQS(); hence this also allows
     more general QSpace tensors X and qvac together with
     the explicit specification of the space index to use.
 
  Wb,Dec28,14
\end{verbatim}

\pagebreak[3]


\myfun{QSpace/getData.m} % l=1/16

\begin{verbatim}
QSpace/getData is a function.
    B = getData(A)
\end{verbatim}

\pagebreak[3]


\myfun{QSpace/getEQdata.m} % l=1/23

\begin{verbatim}
  function [E,Q,D]=getEQdata(HK [,opts])
  Wb,Sep22,12
\end{verbatim}

\pagebreak[3]


\myfun{QSpace/getId.m} % l=6/30

\begin{verbatim}
  function A=getId(A,qs [,k])
 
     reduce QSpace A to rank-2 Id tensor which only has
     vacuum state / scalar representation with each index.
     if k is specified, itag{k} will be used.
 
  Wb,Oct18,14
\end{verbatim}

\pagebreak[3]


\myfun{QSpace/getIdentity.m} % l=5/42

\begin{verbatim}
  function A=getIdentity([args])
 
     same as getIdentityQS
     but returns QSpace object (rather than plain structure).
 
  Wb,May28,16
\end{verbatim}

\pagebreak[3]


\myfun{QSpace/getIdentity3.m} % l=13/53

\begin{verbatim}
  function E=getIdentity3(A,idx [,tag,opts])
 
     Map pair of open indices to be combined as in delta_{ij}
     into single index in E_(i,j,k) by mapping (i=j -> k).
     The added third index therefore always transforms like a scalar
     for all symmetry sectors present, i.e., has q-labels all zero.
 
  Options
 
     --rho normalize E
           such that it resembles density matrix in (i,j) for fixed k.
     tag   tag assigned to the new trailing index.
 
  Wb,Aug28,20
\end{verbatim}

\pagebreak[3]


\myfun{QSpace/getQRange.m} % l=0/13

\begin{verbatim}
  Function qq=getQRange(A,k)
\end{verbatim}

\pagebreak[3]


\myfun{QSpace/getQSpectra.m} % l=17/19

\begin{verbatim}
  function [q,EE,qq,dz,DZ]=getQSpectra(H,ws)
 
    get symmetry-resolved energy spectra.
 
  Input
 
    H  diagonal scalar operator (with diagonals stored only)
    ws specifies which symmetry to pick
 
  Output
 
    q   unique set of labels for symmetry ws
    EE  cell array of energies, one cell for each symmetry in q
    qq  full set of symmetry labels (expanded to same size as EE)
    dz  full set of multiplet degeneracies
    DZ  full set of multiplet degeneracies (expanded to same size as EE)
 
  Wb,Apr12,12
\end{verbatim}

\pagebreak[3]


\myfun{QSpace/get\_cgr\_fac.m} % l=11/42

\begin{verbatim}
  function [dfac,s]=get_cgr_fac(A,i)
 
     CGData is no longer stored with the QSpaces themselves,
     only their weights cgw within outer multiplicity space.
     Therefore CGData is globally normalized to 1. up to
     outer multiplicity.
 
  Options
 
     sflag=='s' => also return range as string info
 
  Wb,Oct17,14
\end{verbatim}

\pagebreak[3]


\myfun{QSpace/get\_qdim.m} % l=4/0

\begin{verbatim}
  function [dd,D]=get_qdim(A,dim)
 
    get cgr dimension for given dimension
 
  Wb,Nov27,15
\end{verbatim}

\pagebreak[3]


\myfun{QSpace/getitags.m} % l=5/10

\begin{verbatim}
  function t=getitags(A)
 
     return itags as cell array.
 
  adapted from QSpace/gotITags.m
  Wb,Apr10,14
\end{verbatim}

\pagebreak[3]


\myfun{QSpace/getproj.m} % l=10/21

\begin{verbatim}
  Function P=getproj(H [,eps,opts])
 
     get projector with respect to H into ground state space
     (diagonalize if necessary); eps determines the window wrt.
     ground state energy, within which state are kept (1E-8).
 
  Options
 
     '-q'  quiet mode
 
  Wb,Aug30,12
\end{verbatim}

\pagebreak[3]


\myfun{QSpace/getqdim.m} % l=4/37

\begin{verbatim}
  function [dd,ss]=getqdim(A)
 
    get symmetry specific length of q-labels
 
  Wb,Sep18,11
\end{verbatim}

\pagebreak[3]


\myfun{QSpace/getqdir.m} % l=9/47

\begin{verbatim}
  function t=getqdir(A [,'-s'])
 
     qdir in numeric format (+- => [+1,-1])
     NB! input arguments are directly handed over to getitags.m
 
  Options
 
     '-s'  return qdir as string (default: numeric; see above)
 
  Wb,Mar30,15
\end{verbatim}

\pagebreak[3]


\myfun{QSpace/getqfmt.m} % l=4/3

\begin{verbatim}
  function f=getqfmt(A [,opts])
 
     get format string for a single Q{i}(j,:)
 
  Wb,Dec14,15
\end{verbatim}

\pagebreak[3]


\myfun{QSpace/getscalar.m} % l=1/13

\begin{verbatim}
  Function a=getscalar(A)
  Wb,Aug28,08
\end{verbatim}

\pagebreak[3]


\myfun{QSpace/getsub.m} % l=33/20

\begin{verbatim}
  Usage #1: function [A,B]=getsub(A,idx)
 
     Select given subspace where idx represents direct index
     into records of QSpace.
 
  Usage #2: function [A,B]=getsub(A,q,dim)
            function [A,B]=getsub(A,{q} [,dim]) // Wb,Feb21,21
 
     Select subspace that has given symmetry sector(s) in given dimension(s).
     if q is a cell array, then this directly relates to to dimensions in A
     e.g. getsub(ops,{[],[],0}) = getsub(ops,0,3) picks scalar sector
     on 3rd dimension; here '0' or 0 can be used as a shortcut to zeros(1,m)
     in q, with m the number of symmetry labels.
 
  Usage #3: function [A,B]=getsub(A,'-0')       // Wb,Dec04,20
            function [A,B]=getsub(A,'-0!')
 
     Assuming A is an irop (rank-3 with irop-index on 3rd leg)
     select scalar operator part (q-labels 0); this is equivalent
     to usage #2, having getsub(A,[0 0 ..],3));
 
     in case of '-0!', also remove singleton irop dimension
     (in case it is not a singleton, i.e., by having multiple
     q=0 multiplets, an error is issued)
 
  Options
 
    -d   return A.data (rather than A itself)
 
    If second output argument is requested, this will contain
    whatever was left in A_in as compared to what was selected
    in A_out.
 
  Wb,Aug20,08 ; Wb,Jun19,13
\end{verbatim}

\pagebreak[3]


\myfun{QSpace/getsym.m} % l=9/0

\begin{verbatim}
  function [ss,s]=getsym(A [,dim,opts])
  Options
 
    -c        return symmetries split as cell array
    -d        get dimensions (rank) of each symmetry [returns 1 for Abelian]
    -r        get symmetry-rank for each symmetry [returns 0 for Abelian]
    'ir',dim  index range in q-labels for given dimension
    '-I',dim  get extended info structure for given symmetry
 
  Wb,Oct24,11
\end{verbatim}

\pagebreak[3]


\myfun{QSpace/getvac.m} % l=11/15

\begin{verbatim}
  function A=getvac(A [,k])
 
     reduce QSpace A to rank-2 Id tensor which only has
     vacuum state / scalar representation with each index.
     if k is specified, itag{k} will be used.
 
  Options
 
     -1d   reduce to rank-1 vector (NB! only permitted for
           for scalar representation with all q-lables zero!)
 
  Wb,Oct18,14
\end{verbatim}

\pagebreak[3]


\myfun{QSpace/getzdim.m} % l=17/32

\begin{verbatim}
  function [dd,qq]=getzdim(A [,dim,opts])
 
    get multiplet dimensions of QSpace A.
 
  Options
 
    -p   product of multiplet dimensions
    -x   expand to match data (repmat)
    -u   get dimension of unique qlabel set (returned as 2nd output argument)
 
   The specification of the dimension (dim) may also take the following
   string values:
 
  '-op'  gets dimension of IROP set
   'op'  gets dimension of state space dimension for scalar operator
         (together with '-u' this is equivalent to 'op') // deprecated
 
  Wb,Jan24,11
\end{verbatim}

\pagebreak[3]


\myfun{QSpace/gotCGS.m} % l=1/55

\begin{verbatim}
  function i = gotCGS(A)
  Wb,May16,10
\end{verbatim}

\pagebreak[3]


\myfun{QSpace/gotITags.m} % l=5/3

\begin{verbatim}
  function r = gotITags(A)
 
     return number of itags (if any).
     second return argument contains listing of ilables as cell array.
 
  Wb,Aug04,12
\end{verbatim}

\pagebreak[3]


\myfun{QSpace/got\_orthogonal\_cgw.m} % l=4/14

\begin{verbatim}
  function got_orthogonal_cgw(A)
 
     (double) check whether A got orthogonal cgw spaces
 
  Wb,Feb19,16
\end{verbatim}

\pagebreak[3]


\myfun{QSpace/hasQOverlap.m} % l=2/24

\begin{verbatim}
  Function: [ss,qq] = hasQOverlap(A,i1,B,i2)
  check overlap in QIDX
  Wb,Jun22,07
\end{verbatim}

\pagebreak[3]


\myfun{QSpace/imag.m} % l=4/32

\begin{verbatim}
  Function: A=imag(A)
 
     extract imaginary data of operator in QSpace format
 
  Wb,Oct19,11
\end{verbatim}

\pagebreak[3]


\myfun{QSpace/info.m} % l=5/42

\begin{verbatim}
  function info(A [,istr,cflag])
 
     header routine used in display()
     Options: istr (info string) and cflag (compact info).
 
  Wb,Mar01,08
\end{verbatim}

\pagebreak[3]


\myfun{QSpace/info2.m} % l=6/53

\begin{verbatim}
  function info2(A,[dim,opts])
  Options
  
    dim  if dimension is specified, plot will be generated
    -ll  log-log scale on plot
  
  show further size info // Wb,Sep17,11
\end{verbatim}

\pagebreak[3]


\myfun{QSpace/info3.m} % l=1/6

\begin{verbatim}
  function info3(A,dim [,opts])
  graphical analysis of size distribution of CGC spaces
\end{verbatim}

\pagebreak[3]


\myfun{QSpace/isAbelian.m} % l=1/13

\begin{verbatim}
  function i=isAbelian(A [,opts])
  Wb,Feb20,13
\end{verbatim}

\pagebreak[3]


\myfun{QSpace/isOpConj.m} % l=1/20

\begin{verbatim}
  function i=isOpConj(A [,output_type])
  Wb,Jun06,19
\end{verbatim}

\pagebreak[3]


\myfun{QSpace/isOrtho.m} % l=1/27

\begin{verbatim}
  function a=isOrtho(A)
  Wb,Jan15,15
\end{verbatim}

\pagebreak[3]


\myfun{QSpace/isQSpace.m} % l=5/34

\begin{verbatim}
  function [i,s]=isQSpace(A [,opts])
 
     overloaded routine (see also MEX/isQSpace.m)
 
  see mpsIsQSpace.m for more information.
  Wb,Feb15,13
\end{verbatim}

\pagebreak[3]


\myfun{QSpace/isdiag.m} % l=18/45

\begin{verbatim}
  function [isd,estr]=isdiag(A)
 
     check whether QSpace is in diagonal rank-2 format.
 
  Return value isd:
 
     1   if block-diagonal in symmetry space
     2   if, in addition, all data{} sets are row-vectors, e.g.
         compressed format as returned bei QSpace::EigenSymmetric()
     3   same as 2, but if all entries in data{} are column-vectors
     0   otherwise
 
  Options
 
   '-d'  If all data blocks are scalars, i.e., of dimension 1x1,
         by default, this returns isd=1; with '-d' this returns
         isd=2, instead.
 
  Wb,Apr24,09
\end{verbatim}

\pagebreak[3]


\myfun{QSpace/isempty.m} % l=1/10

\begin{verbatim}
  function ie=isempty(A)
  Wb,May12,06;  Wb,Aug23,08
\end{verbatim}

\pagebreak[3]


\myfun{QSpace/isequal.m} % l=13/17

\begin{verbatim}
  function i=isequal(A,B [,opts])
  
     overloading isequal and ==
     where input A or B can also be struct( QSpace object ).
 
     Note that the usage A==B is equivalent to isequal(A,B),
     i.e. this assumes the default setting without any options.
 
  Options
 
    '-l'  lenient (ignore difference in itags)
    '-f'  enforce all fields to be the same (no lenience at all)
 
  Wb,Nov09,09 ; Wb,Jun28,21
\end{verbatim}

\pagebreak[3]


\myfun{QSpace/isfield.m} % l=1/36

\begin{verbatim}
  wrapper for builtin isfield
  Wb,Feb26,16
\end{verbatim}

\pagebreak[3]


\myfun{QSpace/isop.m} % l=1/43

\begin{verbatim}
  function is=isop(A)
  Wb,Nov25,11
\end{verbatim}

\pagebreak[3]


\myfun{QSpace/isreal.m} % l=5/50

\begin{verbatim}
  function i=isreal(A)
 
     checks whether A contains complex data
     (this overloads matlabs isreal() function).
 
  Wb,Apr23,15
\end{verbatim}

\pagebreak[3]


\myfun{QSpace/isscalar.m} % l=14/2

\begin{verbatim}
  function i=isscalar(A [,opts])
 
     Checks whether QSpace is a scalar, meaning rank-0
     with a single number in data [e.g., a scalar is obtained
     by tracing/contracting all indices resulting in a single
     number].
 
  Options
 
      -l   lenient (also consider a QSpace with non-empty Q
           a scalar, as long as data contains a single number)
      -x   extended mode (in case of input QSpace array,
           show isscalar() for each individual QSpace
 
  Wb,Aug04,08
\end{verbatim}

\pagebreak[3]


\myfun{QSpace/isscalarop.m} % l=12/22

\begin{verbatim}
  function [s,dd]=isscalarop(A [,opts])
 
     requirements for scalar operators:
      (1) rank-2
      (2) block-diagonal in symmetries
      (3) all cgc spaces are identities
 
  Options
 
    -l   lenient (also accept rank-3 operators)
 
  See also QSpace/issingleop for non-scalar, yet still rank-2 operators.
  Wb,Jul05,12
\end{verbatim}

\pagebreak[3]


\myfun{QSpace/issingleop.m} % l=7/40

\begin{verbatim}
  function [s]=issingleop(A [,opts])
 
     check for possibly non-scalar, yet still rank-2 operators.
 
  See also QSpace/isscalarop for scalar, i.e. block-diagonal
  (rank-2) operators.
 
  Wb,May11,13
\end{verbatim}

\pagebreak[3]


\myfun{QSpace/issingleton.m} % l=1/53

\begin{verbatim}
  Function i=issingleton(A,kk)
  Wb,Apr24,08
\end{verbatim}

\pagebreak[3]


\myfun{QSpace/itag2iter.m} % l=5/1

\begin{verbatim}
  function k=itag2iter(A)
 
     extract site position to k from A.info.itags
     with expected itag format '[non-number]##[non-number]'
 
  Wb,Aug01,17
\end{verbatim}

\pagebreak[3]


\myfun{QSpace/itagrep.m} % l=6/12

\begin{verbatim}
  function A=itagrep(A [,idx], )
 
     Apply given regexprep() pattern to specified itags
     idx specified a subset of itags to consider;
     (otherwise, by default, all itags are considered)
 
  Wb,Jan14,15
\end{verbatim}

\pagebreak[3]


\myfun{QSpace/itags2int.m} % l=9/24

\begin{verbatim}
  function k=itags2int(A)
 
     Extract integer indices carried with itags,
     assuming the itag format '[non-number]##[non-number]'
     where outgoing indices (i.e., which carry a trailing
     conjugate flag `*') are returned as negative integers.
     If no valid number is encountered, nan is returned
     for that index.
 
  Wb,Nov19,21
\end{verbatim}

\pagebreak[3]


\myfun{QSpace/itags2odir.m} % l=1/39

\begin{verbatim}
  function odir=itags2odir(A)
  Wb,May15,17
\end{verbatim}

\pagebreak[3]


\myfun{QSpace/join.m} % l=1/46

\begin{verbatim}
  function A=join(A,idx)
  Wb,May19,10
\end{verbatim}

\pagebreak[3]


\myfun{QSpace/kron.m} % l=0/53

\begin{verbatim}
  overloading kron()
\end{verbatim}

\pagebreak[3]


\myfun{QSpace/logical.m} % l=6/0

\begin{verbatim}
  function q=logical(A)
 
     Overload of conversion to logicals (unary operator)
     e.g. for conditionals: if A, ...; else ...; end
 
  see also QSpace/ [which returns ~logical()]
  Wb,Feb18,21
\end{verbatim}

\pagebreak[3]


\myfun{QSpace/makeIrop.m} % l=13/12

\begin{verbatim}
  function A=makeIrop(A [,t_op])
 
     make A an irreducible operator with a proper 3rd irop index,
     which thus is only intended for rank-2 scalar operators.
     Operators already in rank-3 format are checked for the
     right format.
 
  Options
 
     t_op specifies itag for the irop leg, only if added
     (empty by default).
 
  See also: fixScalarOp, appendSingletons.
  Wb,Feb12,16
\end{verbatim}

\pagebreak[3]


\myfun{QSpace/mapQ.m} % l=8/31

\begin{verbatim}
  QM = Q map that transforms Q numbers to `rotated' set
  Example: SIAM: (N1,N2) --> (Q,Sz)
 
       QM      *  Q_old    + offset  =  Q_new
 
     [ 1  1    *  [ N1     +  [-1    =  [ Q
       1 -1 ]       N2 ]        0 ]       Sz ]
 
  Wb,Aug10,6
\end{verbatim}

\pagebreak[3]


\myfun{QSpace/minus.m} % l=0/45

\begin{verbatim}
  overloading the + operator
\end{verbatim}

\pagebreak[3]


\myfun{QSpace/mrdivide.m} % l=1/51

\begin{verbatim}
  overloading * operator
  Wb,Sep28,12
\end{verbatim}

\pagebreak[3]


\myfun{QSpace/mtimes.m} % l=1/58

\begin{verbatim}
  function C=mtimes(A,B,ia,ib)
  overloading * operator (matrix multiply)
\end{verbatim}

\pagebreak[3]


\myfun{QSpace/norm.m} % l=0/6

\begin{verbatim}
  wrapper routine for normQS()
\end{verbatim}

\pagebreak[3]


\myfun{QSpace/norm2.m} % l=1/12

\begin{verbatim}
  wrapper routine norm2 -> normQS()
  Wb,Oct18,14
\end{verbatim}

\pagebreak[3]


\myfun{QSpace/not.m} % l=5/19

\begin{verbatim}
  function q=not(A)
 
     Overload of conversion to logicals (unary operator)
     e.g. for conditionals: if ~A, ...; else ...; end
 
  Wb,Feb18,21
\end{verbatim}

\pagebreak[3]


\myfun{QSpace/numsym.m} % l=4/30

\begin{verbatim}
  function ns=numsym(A)
 
     determine number of symmetries in QSpace(s) A.
 
  Wb,Feb06,14
\end{verbatim}

\pagebreak[3]


\myfun{QSpace/olap.m} % l=4/40

\begin{verbatim}
  function x=olap(A [,B,opts])
 
     calculate frobenius norm / overlap.
 
  Wb,May27,11
\end{verbatim}

\pagebreak[3]


\myfun{QSpace/oplus.m} % l=15/50

\begin{verbatim}
  function C=oplus(A,B [r12,opts]
 
     Direct sum (latex oplus) of two tensors A and B along dimension r.
     If r12 is not specified, an MPS or MPO like setting is assumed
     with L/R as first and second index, i.e., r12=[1 2].
 
  Options
 
    'first'             direct sum only in index r12(2)
    'last'              direct sum only in index r12(1)
                        otherwise: direct sum in both indices r12
 
    'bfac',..           apply bfac to B prior to extending A
 
  E.g. see https://mathworld.wolfram.com/MatrixDirectSum.html
  Wb,Nov23,20
\end{verbatim}

\pagebreak[3]


\myfun{QSpace/permute.m} % l=2/12

\begin{verbatim}
  Usage: A=permute(A)
  permute QSpace with given permutation
  Wb,Sep08,06 ; Wb,Nov24,14
\end{verbatim}

\pagebreak[3]


\myfun{QSpace/plotEdata.m} % l=6/20

\begin{verbatim}
  function plotEdata(H [,opts])
 
     plot eigenspectrum of given QSpace,
     expecting scalar operator upon input!
 
  see also QSpace/
  Wb,Feb15,13
\end{verbatim}

\pagebreak[3]


\myfun{QSpace/plotQSpectra.m} % l=24/32

\begin{verbatim}
  function Iq=plotQSpectra([H|R] [,opts][,SOP])
 
     graphical analysis of eigenspectrum of Hamiltonian
     (or entanglement spectrum of density matrix; requires '-ES' flag)
     separating the data into its symmetry subspaces.
 
  Options
 
    '-ES'     entanglement spectrum (expects Rho as input)
    '-x'      show state space dimensions (rather than multiplet dimensions)
 
    'yl',...  ylim range
    'y2',...  y-range within which degeneracies are indicated
    'ws',...  which symmetries to pick (e.g. [1 2] for the first two)
 
    SOP       more verbose description of symmetry labels, which also
              appears in the title of the respective panels.
              SOP my simply be specified through IS.SOP, i.e. as returned
              by getLocalSpace() in the NRG setup; all that is used, however,
              are the strings SOP(:).info, hence, alternatively, SOP may
              only contain {SOP.info}. keywords looked for are 'spin' or
              'charge' which effects the xlabels.
 
  Adapted from getrhoESpectra.m
  Wb,Apr11,12
\end{verbatim}

\pagebreak[3]


\myfun{QSpace/plot\_op.m} % l=5/3

\begin{verbatim}
  function plot_op(S,H)
 
     S operator (possibly IROP)
     H Hamiltonian (in diagonal representation) => show energy range
 
  Wb,Aug20,12
\end{verbatim}

\pagebreak[3]


\myfun{QSpace/plus.m} % l=0/14

\begin{verbatim}
  overloading the + operator
\end{verbatim}

\pagebreak[3]


\myfun{QSpace/randomize.m} % l=9/20

\begin{verbatim}
  function A=randomize(A [,nrm,opts])
 
     Randomize A.data using randn()
     while preserving overall norm (unless explicitly specified).
 
  Options
 
     -z   use complex random numbers (default: real)
 
  Wb,Nov11,09
\end{verbatim}

\pagebreak[3]


\myfun{QSpace/rank.m} % l=0/35

\begin{verbatim}
  overloading rank function
\end{verbatim}

\pagebreak[3]


\myfun{QSpace/real.m} % l=4/41

\begin{verbatim}
  Function: A=real(A)
 
     extract real data of operator in QSpace format
 
  Wb,Oct19,11
\end{verbatim}

\pagebreak[3]


\myfun{QSpace/rgbQSpectra.m} % l=13/51

\begin{verbatim}
  function [cc,Iout]=rgbQSpectra(H,q3 [,opts])
 
    Encode given energy spectrum to rgb specturm
    by using thermal weights in specified (or otherwise dominant)
    symmetry sectors
 
  Options
 
    'beta',.. effective inverse `temperature' to use to compute
             `thermal' weight in symmetry sectors (2)
    '-R'      assume input is density matrix which already gives
              weights (with beta ignored here)
 
  Wb,Sep09,21
\end{verbatim}

\pagebreak[3]


\myfun{QSpace/rmitags.m} % l=5/11

\begin{verbatim}
  function A=rmitags(A)
 
     skip itags of given QSpace(s)
 
  *** NB! deprecated; rather use untag ***
  Wb,Apr14,16
\end{verbatim}

\pagebreak[3]


\myfun{QSpace/rmtrace.m} % l=4/22

\begin{verbatim}
  function A=rmtrace(A)
 
     remove trace from input tensor A
 
  Wb,Apr24,13
\end{verbatim}

\pagebreak[3]


\myfun{QSpace/sameas.m} % l=6/32

\begin{verbatim}
  function i=sameas(A,B [,eps=1E-10])
 
     check whether QSpaces A and B are the same
     except for minor differences such as 
     info.cgr.cid(end) and info.cgr.nnz
 
  Wb,Aug08,16
\end{verbatim}

\pagebreak[3]


\myfun{QSpace/setitags.m} % l=41/44

\begin{verbatim}
  function A=setitags(A,...)
 
    set itag(s) while preserving conj flag.
 
  Usage #1: A=setitags(A,{t1,t2,...} [,[i1,i2,...]])
 
    set specified or all itags
    setitags(A,{'K1','K2','s2'})        sets itags { 'K1','K2,'s2' }
    setitags(A,{'K','K','s'},{1,2,2})   same (numbers are appended to itags!)
 
  Usage #2: A=setitags(A,tt,i)
 
    set specific itags tt to itag{i}
 
  Usage #3: set itags for A-tensor (conj-flags always remain in place)
 
    A=setitags(A,'-A',k)           {'K<k-1>','K<k>','s<k>}
    A=setitags(A,'-A:K,K,s',  k)   {'K<k-1>','K<k>','s<k>} more specific
    A=setitags(A,'-A>K,K,s@n',k)   same, also using n-digit format (default: 2)
    A=setitags(A,'-A<K,K,s@n',k)   {'K<k>','K<k+1>','s<k>} i.e. R->L order
 
  Usage #4: set itags for operator
 
    S=setitags(S,'-op:X', k )       {'<X><k>','<X><k>','op'}
    S=setitags(S,'-op:X','s')       {'<X>','<X>','s'}
    S=setitags(S,'-op:<regex>', A)  set itags based on behavior
                                    that mimicks contractQS
  Usage #5
 
    A=setitags(A,ia,B,ib);  copy itags from B to A
 
  Usage #6: any of the usages above, but return itags rather than A,
  by requesting 2 output args, e.g.
 
    [tt,i]=setitags(A,...)
    [tt,~]=setitags(A,...)
 
    this returns itags+indizes themselves, rather than setting the
    itags in A (the conjugation is inherited from A; hence only a
    single QSpace A is expected).
 
  Wb,Mar24,16 ; Wb,May28,18
\end{verbatim}

\pagebreak[3]


\myfun{QSpace/showdeg.m} % l=13/32

\begin{verbatim}
  function [I=]showdeg(H [,opts])
 
     show degeneracy of given rank-2 operator, either in 
     diagonal representation, or to be diagonalized.
 
  Options
 
    'dE',...  considere values degenerate within dE (1E-3)
    'ng',...  number of degenerate groups to display
    '-f'      force flag
    '-v'      always print lowest part of spectrum
    '-V'      more verbose still.
 
  Wb,Apr12,11
\end{verbatim}

\pagebreak[3]


\myfun{QSpace/skiprecs.m} % l=1/51

\begin{verbatim}
  function A=skiprecs(A,ix)
  Wb,Aug01,12
\end{verbatim}

\pagebreak[3]


\myfun{QSpace/skipzeros.m} % l=1/58

\begin{verbatim}
  function A=skipzeros(A [,eps,'--all'])
  Wb,Nov10,07
\end{verbatim}

\pagebreak[3]


\myfun{QSpace/sort.m} % l=10/6

\begin{verbatim}
  function [A,isd]=sort(A,perm)
 
     sort records in QSpace according to q-labels after using
     permutation perm (default: Id). Alternatively, perm may
     also be specified as 
 
       '-E' (sort by energy) or
       '-R' (same as '-E', but sorts in descending order).
       '-S' sort wrt. to size of data+cgr
 
  Wb,Sep12,06
\end{verbatim}

\pagebreak[3]


\myfun{QSpace/split\_ops.m} % l=15/22

\begin{verbatim}
  function [A,nsp]=split_ops(A [,opts])
 
     split composite operators back to (irreducible) operator sets
     that have well defined irop symmetry labels
     (returns A -> A(:) expanded along dim2, such that
     A(k) still references a (partial) original QSpace).
 
  Options
 
     --set-op    set as operator type
     --s2f       move scalar operators to front
     --trans     return transpose array
     --vec       return as serialzed vector
     --cell      return as cell array
 
  Wb,Aug31,19
\end{verbatim}

\pagebreak[3]


\myfun{QSpace/sqrt.m} % l=5/43

\begin{verbatim}
  function B=sqrt(A)
 
     Compute sqrt(A) for diagonal QSpace(s) A
     (error is issued if A is not diagonal)
 
  Wb,Mar11,19
\end{verbatim}

\pagebreak[3]


\myfun{QSpace/squeeze.m} % l=9/54

\begin{verbatim}
  Function [A,I]=squeeze(A [,OPTS])
 
     remove singleton dimensions (similar to matlab's routine)
 
  Options
 
    '-start'  only singletons at the front
    '-end'    only singletons at the end (default: all)
 
  Wb,Apr24,08 ; Wb,Dec01,20
\end{verbatim}

\pagebreak[3]


\myfun{QSpace/subsasgn.m} % l=4/10

\begin{verbatim}
  subsasgn
 
  Reference http://www.cs.ubc.ca/~murphyk/Software/matlabObjects.html
  Classes and objects in matlab: the quick and dirty way
  Kevin Murphy, 2005
\end{verbatim}

\pagebreak[3]


\myfun{QSpace/subsref.m} % l=4/20

\begin{verbatim}
  subsref
 
  Reference http://www.cs.ubc.ca/~murphyk/Software/matlabObjects.html
  Classes and objects in matlab: the quick and dirty way
  Kevin Murphy, 2005
\end{verbatim}

\pagebreak[3]


\myfun{QSpace/sum.m} % l=0/30

\begin{verbatim}
  overloading the sum operator
\end{verbatim}

\pagebreak[3]


\myfun{QSpace/symperm.m} % l=4/36

\begin{verbatim}
  function A=symperm(A,perm)
 
     Permute symmetries according to specified permutation.
 
  Wb,Feb06,14
\end{verbatim}

\pagebreak[3]


\myfun{QSpace/symrank.m} % l=1/46

\begin{verbatim}
  function r=symrank(A)
  Wb,Jan09,17
\end{verbatim}

\pagebreak[3]


\myfun{QSpace/tensor.m} % l=2/53

\begin{verbatim}
  function C=tensor(A,B,...)
  kron including merging of quantum numbers
  Wb,Jul07,09
\end{verbatim}

\pagebreak[3]


\myfun{QSpace/times.m} % l=6/2

\begin{verbatim}
  overloading .* operator
  Options
 
    --skip-cgc   C inherits CGC space of tensor A *as it is*
                 (yet only for the symmetry sectors shared with B).
 
  Wb,Jul04,12
\end{verbatim}

\pagebreak[3]


\myfun{QSpace/toStr.m} % l=1/14

\begin{verbatim}
QSpace/toStr is a function.
    qs = toStr(A, varargin)
\end{verbatim}

\pagebreak[3]


\myfun{QSpace/trace.m} % l=10/21

\begin{verbatim}
  Function x=trace(A,I [,opts])
 
    trace of QSpace operator
 
  Usage 1: A=trace(A)   - regular trace of rank-2 object
  Usage 2: A=trace(A,I)
 
     generalized trace with pairwise specification of indizes
     to contract I = [i1 i2; j1 j2; ... ]
 
  Wb,Sep11,06 ; Wb,Apr10,15
\end{verbatim}

\pagebreak[3]


\myfun{QSpace/transpose.m} % l=0/37

\begin{verbatim}
  overloading the .' operator (element wise transpose)
\end{verbatim}

\pagebreak[3]


\myfun{QSpace/tstfun.m} % l=1/43

\begin{verbatim}
  function tstfun(A)
  Wb,Feb23,11
\end{verbatim}

\pagebreak[3]


\myfun{QSpace/uminus.m} % l=4/50

\begin{verbatim}
  function A=uminus(A)
 
     overloading the unary minus operator (as in -A)
 
  Wb,Sep17,06
\end{verbatim}

\pagebreak[3]


\myfun{QSpace/untag.m} % l=7/1

\begin{verbatim}
  Usage #1: function A=untag(A)
  Usage #2: function untag(A1,A2,...)
 
     remove itags from specified A's;
     if no return argument is requested, the modified A's
     are assigned back to caller by their inputname.
 
  Wb,May11,17
\end{verbatim}

\pagebreak[3]


\myfun{QSpace/uplus.m} % l=4/14

\begin{verbatim}
  function A=uplus(A)
 
     overloading the unary plus operator (as in +A)
 
  Wb,Sep17,06
\end{verbatim}

\end{document}
